\documentclass[12pt,a4paper]{article}

\usepackage[utf8]{inputenc}
\usepackage[T1]{fontenc}
\usepackage{amsmath,amssymb,amsthm}
\usepackage{mathtools}
\usepackage{bm}
\usepackage{geometry}
\usepackage{hyperref}
\usepackage{xcolor}
\usepackage{tcolorbox}
\usepackage{booktabs}
\usepackage{enumitem}
\usepackage{algorithm}
\usepackage{algpseudocode}
\usepackage{tikz}
\usepackage{pgfplots}
\usepackage{mdframed}
\usepackage{fancyhdr}
\usepackage{titlesec}
\usepackage{parskip}
\usepackage{microtype}
\usepackage{array}
\usepackage{longtable}
\usepackage{multirow}
\usepackage{makecell}

\geometry{
  left=2.5cm,
  right=2.5cm,
  top=2.5cm,
  bottom=2.5cm
}

\pgfplotsset{compat=1.18}

\hypersetup{
  colorlinks=true,
  linkcolor=blue!70!black,
  citecolor=green!50!black,
  urlcolor=blue
}

% Color definitions
\definecolor{defcolor}{RGB}{0,70,127}
\definecolor{thmcolor}{RGB}{127,0,0}
\definecolor{notecolor}{RGB}{0,100,0}
\definecolor{keycolor}{RGB}{90,0,90}
\definecolor{excolor}{RGB}{180,90,0}
\definecolor{alertcolor}{RGB}{200,0,0}
\definecolor{boxbg}{RGB}{240,245,255}
\definecolor{notebg}{RGB}{240,255,240}
\definecolor{alertbg}{RGB}{255,240,240}
\definecolor{keybg}{RGB}{248,240,255}
\definecolor{headcolor}{RGB}{0,50,120}

% Theorem environments
\tcbuselibrary{theorems,skins,breakable}

\newtcolorbox{definitionbox}[1][]{
  enhanced,
  colback=boxbg,
  colframe=defcolor,
  fonttitle=\bfseries\color{white},
  coltitle=white,
  attach boxed title to top left={yshift=-2mm, xshift=4mm},
  boxed title style={colback=defcolor, rounded corners},
  title={Definition},
  breakable,
  #1
}

\newtcolorbox{theorembox}[1][]{
  enhanced,
  colback=alertbg,
  colframe=thmcolor,
  fonttitle=\bfseries\color{white},
  coltitle=white,
  attach boxed title to top left={yshift=-2mm, xshift=4mm},
  boxed title style={colback=thmcolor, rounded corners},
  title={Theorem / Key Result},
  breakable,
  #1
}

\newtcolorbox{notebox}[1][]{
  enhanced,
  colback=notebg,
  colframe=notecolor,
  fonttitle=\bfseries\color{white},
  coltitle=white,
  attach boxed title to top left={yshift=-2mm, xshift=4mm},
  boxed title style={colback=notecolor, rounded corners},
  title={Note},
  breakable,
  #1
}

\newtcolorbox{keybox}[1][]{
  enhanced,
  colback=keybg,
  colframe=keycolor,
  fonttitle=\bfseries\color{white},
  coltitle=white,
  attach boxed title to top left={yshift=-2mm, xshift=4mm},
  boxed title style={colback=keycolor, rounded corners},
  title={Key Insight},
  breakable,
  #1
}

\newtcolorbox{examplebox}[1][]{
  enhanced,
  colback=yellow!5,
  colframe=excolor,
  fonttitle=\bfseries\color{white},
  coltitle=white,
  attach boxed title to top left={yshift=-2mm, xshift=4mm},
  boxed title style={colback=excolor, rounded corners},
  title={Example},
  breakable,
  #1
}

\newtcolorbox{alertbox}[1][]{
  enhanced,
  colback=alertbg,
  colframe=alertcolor,
  fonttitle=\bfseries\color{white},
  coltitle=white,
  attach boxed title to top left={yshift=-2mm, xshift=4mm},
  boxed title style={colback=alertcolor, rounded corners},
  title={Important},
  breakable,
  #1
}

% Header and footer
\pagestyle{fancy}
\fancyhf{}
\fancyhead[L]{\color{headcolor}\textbf{Autoencoders, Regularisation \& Beyond}}
\fancyhead[R]{\color{headcolor}\textit{IIIT Hyderabad -- Lecture Notes}}
\fancyfoot[C]{\thepage}
\renewcommand{\headrulewidth}{0.5pt}
\renewcommand{\footrulewidth}{0.3pt}

% Section formatting
\titleformat{\section}[block]
  {\Large\bfseries\color{headcolor}}
  {\thesection}{1em}{}[\titlerule]

\titleformat{\subsection}[block]
  {\large\bfseries\color{defcolor}}
  {\thesubsection}{1em}{}

\titleformat{\subsubsection}[block]
  {\normalsize\bfseries\color{keycolor}}
  {\thesubsubsection}{1em}{}

% Math macros
\newcommand{\R}{\mathbb{R}}
\newcommand{\E}{\mathbb{E}}
\newcommand{\Var}{\mathrm{Var}}
\newcommand{\Bias}{\mathrm{Bias}}
\newcommand{\Cov}{\mathrm{Cov}}
\newcommand{\sigmoid}{\sigma}
\newcommand{\norm}[1]{\left\|#1\right\|}
\newcommand{\abs}[1]{\left|#1\right|}
\newcommand{\inner}[2]{\langle #1, #2 \rangle}
\newcommand{\grad}{\nabla}
\newcommand{\vx}{\mathbf{x}}
\newcommand{\vh}{\mathbf{h}}
\newcommand{\vw}{\mathbf{w}}
\newcommand{\vb}{\mathbf{b}}
\newcommand{\vc}{\mathbf{c}}
\newcommand{\vz}{\mathbf{z}}
\newcommand{\hatx}{\hat{\mathbf{x}}}
\newcommand{\Loss}{\mathcal{L}}
\newcommand{\Reg}{\mathcal{S}}

%=====================================================================
\begin{document}
%=====================================================================

\begin{titlepage}
  \centering
  \vspace*{2cm}
  {\Huge\bfseries\color{headcolor}
    Autoencoders, Regularisation,\\[0.5em]
    Bias--Variance and Beyond\\[1em]}
  {\Large\color{defcolor} Comprehensive Lecture Notes with Full Derivations}\\[2em]
  \rule{\textwidth}{1.5pt}\\[1em]
  {\large\textit{Course: Machine Learning / Deep Learning}}\\[0.5em]
  {\large\textit{International Institute of Information Technology, Hyderabad}}\\[0.5em]
  {\large\textit{Under the Guidance of Prof.\ Praveen Parchuri}}\\[2em]
  \rule{\textwidth}{0.5pt}\\[1em]
  {\normalsize Covers: Autoencoders $\cdot$ PCA Equivalence $\cdot$ Loss Functions $\cdot$ Backpropagation\\
  Regularisation (L2, Weight Tying, Denoising, Sparse, Contractive) $\cdot$
  Bias-Variance Tradeoff\\
  Training vs.\ Test Error $\cdot$ Stein's Lemma $\cdot$ Dataset Augmentation $\cdot$ Early Stopping\\
  Ensemble Methods $\cdot$ Dropout $\cdot$ Discrete Cosine Transform $\cdot$ JPEG Compression}
  \vfill
  {\large\today}
\end{titlepage}

\tableofcontents
\newpage

%=====================================================================
\section{Introduction to Autoencoders}
%=====================================================================

\subsection{What Is an Autoencoder?}

\begin{definitionbox}
An \textbf{autoencoder} is a special type of feed-forward neural network that learns to:
\begin{enumerate}
  \item \textbf{Encode} its input $\vx_i \in \R^n$ into a hidden (latent) representation
    $\vh = g(W\vx_i + \vb) \in \R^d$
  \item \textbf{Decode} that hidden representation back to a reconstructed output
    $\hatx_i = f(W^*\vh + \vc) \in \R^n$
\end{enumerate}

The architecture is summarised as:
\[
  \vx_i
    \;\xrightarrow{\;W,\,\vb\;}\;
  \vh
    \;\xrightarrow{\;W^*,\,\vc\;}\;
  \hatx_i
\]
\end{definitionbox}

\paragraph{Notation summary:}
\begin{itemize}
  \item $\vx_i \in \R^n$: the $i$-th input data point (data matrix $X \in \R^{m\times n}$, $m$ samples, $n$ features)
  \item $\vh \in \R^d$: the hidden / latent representation; $\vh = g(W\vx_i + \vb)$
  \item $\hatx_i \in \R^n$: reconstructed output; $\hatx_i = f(W^*\vh + \vc)$
  \item $g(\cdot)$: encoder nonlinearity (e.g.\ sigmoid, tanh, ReLU)
  \item $f(\cdot)$: decoder nonlinearity (depends on input type)
  \item $W \in \R^{d \times n}$: encoder weight matrix
  \item $W^* \in \R^{n \times d}$: decoder weight matrix (note the transpose dimensions)
  \item $\vb \in \R^d$: encoder bias; $\vc \in \R^n$: decoder bias
  \item Parameter space dimension: $R^d = R^{d\times n} \cdot R^{n+1} \cdot R^{d+1}$
\end{itemize}

\begin{keybox}
\textbf{Goal:} Learn what the \emph{most important characteristics} of $\vx_i$ are, such that $\vh$ captures everything needed to reconstruct the original input with minimal loss.
\end{keybox}

%-------------------------------------------------------------------
\subsection{The Compression Perspective and Analogy with PCA}
%-------------------------------------------------------------------

Consider the case where $\dim(\vh) < \dim(\vx_i)$, i.e.\ $d < n$. If we can \emph{still} reconstruct $\hatx_i$ from $\vh$, then $\vh$ is a \textbf{loss-free (or low-loss) encoding}. It must capture all the important characteristics of $\vx_i$.

\begin{examplebox}
Think of it like file compression: $10\text{ bits} \to 4\text{ bits} \to 10\text{ bits}$. If the decoded file equals the original, the 4-bit representation captured everything essential.
\end{examplebox}

\paragraph{Analogy with PCA:}
\begin{itemize}
  \item \textbf{PCA}: Projects data onto the top-$k$ principal components, discarding noise / low-variance / correlated dimensions.
  \item \textbf{Autoencoder hidden layer $\vh$}: Performs similar dimensionality reduction, but \emph{nonlinearly}, allowing it to capture more complex structure.
  \item Both seek a \textbf{compact representation} of the data.
\end{itemize}

%-------------------------------------------------------------------
\subsection{Under-Complete vs.\ Over-Complete Autoencoders}
%-------------------------------------------------------------------

\begin{definitionbox}
\begin{itemize}
  \item \textbf{Under-complete autoencoder}: $d < n$ (hidden dimension smaller than input). Forces the network to learn a concise summary -- it \emph{cannot} trivially copy input.
  \item \textbf{Over-complete autoencoder}: $d \geq n$ (hidden dimension at least as large as input). The network \emph{could} learn a trivial identity: copy $\vx_i$ into $\vh$ and then copy $\vh$ into $\hatx_i$.
\end{itemize}
\end{definitionbox}

\begin{alertbox}
\textbf{Key Problem with Over-Complete AEs:} Without additional constraints, an over-complete autoencoder has \emph{no incentive} to learn a meaningful representation. It will simply memorise the training data. This motivates \textbf{regularisation}.
\end{alertbox}

\begin{notebox}
\textbf{BMI Example:} Even if you want to extract $(\text{height}, \text{weight}, \text{BMI})$ from data, since $\text{BMI} = f(\text{height}, \text{weight})$ is deterministic, there is no reason why the machine would disentangle BMI without special methods. An over-complete AE would simply copy the inputs.
\end{notebox}

The training objective is to minimise a \textbf{reconstruction loss} that ensures $\hatx_i$ is close to $\vx_i$.

%=====================================================================
\section{Loss Functions for Autoencoders}
%=====================================================================

\subsection{Loss for Real-Valued Inputs: Mean Squared Error}

When inputs $\vx_i \in \R^n$ are real-valued, the natural loss is the \textbf{mean squared error (MSE)}:

\[
  \Loss(\theta) = \frac{1}{m}\sum_{i=1}^{m}\sum_{j=1}^{n}
    \bigl(\hat{x}_{ij} - x_{ij}\bigr)^2
\]

In matrix notation:
\[
  \Loss(\theta) = \frac{1}{m}\sum_{i=1}^{m}
    (\vx_i - \hatx_i)^\top(\hatx_i - \vx_i)
\]

\paragraph{Decoder choice for real-valued inputs:}
\begin{itemize}
  \item Sigmoid and tanh restrict outputs to $[0,1]$ or $[-1,1]$, but $\vx_i \in \R^n$.
  \item Most appropriate decoder: \textbf{linear}, i.e.\ $\hatx_i = W^*\vh + \vc \in \R^n$.
  \item The encoder $g$ is typically chosen as the \textbf{sigmoid} function.
\end{itemize}

\begin{examplebox}
$\vx_i = [0.25,\; 0.5,\; 0.5]^\top$ with real-valued inputs. The decoder should be linear (not restricted to $[0,1]$).
\end{examplebox}

%-------------------------------------------------------------------
\subsection{Loss for Binary Inputs: Cross-Entropy}
%-------------------------------------------------------------------

When inputs $x_{ij} \in \{0, 1\}$ are binary, a probabilistic interpretation is natural.

\paragraph{Decoder choice:} Use a \textbf{sigmoid (logistic)} decoder:
\[
  \hat{x}_{ij} = \sigma(a_j) = \frac{1}{1 + e^{-a_j}}
\]
where $a_j$ is the $j$-th pre-activation of the output layer. The output is interpreted as the probability that bit $j$ is 1:
\begin{itemize}
  \item Output $0.8$ means ``80\% probability that this bit is 1''
  \item Output $0.2$ means ``20\% probability that this bit is 1 (80\% probability of 0)''
\end{itemize}

\paragraph{Cross-entropy loss (preferred over squared loss for binary inputs):}
\[
  \Loss(\theta) = -\frac{1}{m}\sum_{i=1}^{m}\sum_{j=1}^{n}
    \Bigl[x_{ij}\log\hat{x}_{ij} + (1-x_{ij})\log(1-\hat{x}_{ij})\Bigr]
\]

%-------------------------------------------------------------------
\subsection{Deriving Cross-Entropy from First Principles}
%-------------------------------------------------------------------

\paragraph{Probabilistic setup:}
Let $p = [x_{ij},\; 1-x_{ij}]^\top$ be the true binary distribution and $q = [\hat{x}_{ij},\; 1-\hat{x}_{ij}]^\top$ be the estimated distribution. The \textbf{cross-entropy} between $p$ and $q$ is:
\[
  H(p,q) = -\sum_{k=0}^{1} p_k \log q_k
\]

Applying this to each data point $i$ and feature $j$:
\begin{align*}
  \ell_{ij}
  &= -\bigl[x_{ij}\log\hat{x}_{ij} + (1-x_{ij})\log(1-\hat{x}_{ij})\bigr]\\
  &= \begin{cases}
      -\log\hat{x}_{ij} & \text{if } x_{ij}=1\\
      -\log(1-\hat{x}_{ij}) & \text{if } x_{ij}=0
     \end{cases}
\end{align*}

This is minimised when $\hat{x}_{ij} = x_{ij}$, i.e.\ $\hat{x}_{ij}=1$ when $x_{ij}=1$ and $\hat{x}_{ij}=0$ when $x_{ij}=0$.

\begin{keybox}
\textbf{Why Cross-Entropy $\succ$ Squared Loss for Binary?}\\
Cross-entropy naturally aligns with the Bernoulli probabilistic model. It avoids saturation in gradient computation (the sigmoid gradient cancels out cleanly, see backpropagation section) and works much better in practice for binary outputs.
\end{keybox}

%-------------------------------------------------------------------
\subsection{Autoencoder Architecture Summary}
%-------------------------------------------------------------------

\begin{center}
\begin{tabular}{lp{5cm}p{5cm}}
\toprule
 & \textbf{Binary Inputs} & \textbf{Real-Valued Inputs} \\
\midrule
Encoder $g$ & Sigmoid & Sigmoid or tanh \\
Decoder $f$ & Sigmoid & Linear ($W^*\vh + \vc$) \\
Loss & Cross-entropy & Mean Squared Error \\
\bottomrule
\end{tabular}
\end{center}

The hidden layer with dimension $d$ is the \textbf{bottleneck}. When $d < n$ (under-complete), the model is forced to compress. The bottleneck dimension controls capacity.

%=====================================================================
\section{Backpropagation Through an Autoencoder}
%=====================================================================

\subsection{Training Like a Regular Network}

An autoencoder is trained exactly like a feed-forward neural network using \textbf{backpropagation}. We need:
\[
  \frac{\partial\Loss(\theta)}{\partial W^*},\quad
  \frac{\partial\Loss(\theta)}{\partial W},\quad
  \frac{\partial\Loss(\theta)}{\partial \vc},\quad
  \frac{\partial\Loss(\theta)}{\partial \vb}
\]

The weight $W$ appears in \emph{two} places:
\begin{enumerate}
  \item The encoder path: $\vh = g(W\vx_i + \vb)$
  \item The decoder path: $\hatx_i = f(W^*\vh + \vc)$
\end{enumerate}

Therefore:
\[
  \frac{\partial\Loss}{\partial W}
  = \left(\frac{\partial\Loss}{\partial W}\right)_{\text{encoder}}
  + \left(\frac{\partial\Loss}{\partial W}\right)_{\text{decoder}}
\]

%-------------------------------------------------------------------
\subsection{Backprop: Decoder Contribution (Squared Loss)}
%-------------------------------------------------------------------

For the squared-error loss $\Loss = (\hatx_i - \vx_i)^\top(\hatx_i - \vx_i)$ with linear decoder $\hatx_i = f(W^*\vh + \vc) = W^*\vh + \vc$:

\begin{align*}
  \frac{\partial\Loss}{\partial\hatx_i}
  &= 2(\hatx_i - \vx_i)
\end{align*}

For the \textbf{decoder weights} $W^*$, let $\vz_d = W^*\vh + \vc$ (pre-activation of output layer):
\[
  \frac{\partial\Loss}{\partial W^*}
  = \frac{\partial\Loss}{\partial \vz_d} \cdot \frac{\partial \vz_d}{\partial W^*}
  = \delta_d \cdot \vh^\top
\]
where $\delta_d = \frac{\partial\Loss}{\partial \vz_d} = 2(\hatx_i - \vx_i)$ is the output-layer error signal (for linear decoder).

%-------------------------------------------------------------------
\subsection{Backprop: Encoder Contribution}
%-------------------------------------------------------------------

For the encoder weights $W$, gradients flow through both the decoder and encoder via the chain rule:
\[
  \frac{\partial\Loss}{\partial W}
  = \underbrace{\frac{\partial\Loss}{\partial\hatx_i}}_{\text{reused}}
    \cdot \underbrace{\frac{\partial\hatx_i}{\partial\vh}}_{\text{reused}}
    \cdot \underbrace{\frac{\partial\vh}{\partial(W\vx_i+\vb)}}_{\text{recomputed}}
    \cdot \frac{\partial(W\vx_i+\vb)}{\partial W}
\]

In practice, for a sigmoid encoder $\vh = \sigma(W\vx_i + \vb)$:
\[
  \frac{\partial\Loss}{\partial W}
  = \delta_e \cdot \vx_i^\top
\]
where $\delta_e$ is propagated error at the encoder via backpropagation.

\paragraph{Key gradient for binary cross-entropy + sigmoid output:}
With $h_{2j} = \hat{x}_{ij}$ and cross-entropy loss:
\begin{align*}
  \frac{\partial\Loss}{\partial h_{2j}}
  &= -\frac{x_{ij}}{\hat{x}_{ij}} + \frac{1-x_{ij}}{1-\hat{x}_{ij}}
\end{align*}

Sigmoid derivative:
\[
  \frac{dh_{2j}}{da_{2j}} = \sigma(a_{2j})(1 - \sigma(a_{2j}))
\]

The beautiful simplification (sigmoid + cross-entropy):
\[
  \frac{\partial\Loss}{\partial a_{2j}}
  = \frac{\partial\Loss}{\partial h_{2j}} \cdot \frac{dh_{2j}}{da_{2j}}
  = \left(-\frac{x_{ij}}{\hat{x}_{ij}} + \frac{1-x_{ij}}{1-\hat{x}_{ij}}\right) \cdot \hat{x}_{ij}(1-\hat{x}_{ij})
  = \hat{x}_{ij} - x_{ij}
\]

This is remarkably clean -- the gradient is simply the prediction error.

\begin{keybox}
\textbf{Key observation:} The gradients always have the same \emph{shape} as the variable you're differentiating with respect to. This is why backpropagation is so elegant.
\end{keybox}

Full chain for encoder weights:
\[
  \frac{\partial\Loss}{\partial W}
  = \frac{\partial\Loss}{\partial\vh_2}
    \cdot \frac{\partial\vh_2}{\partial a_2}
    \cdot \frac{\partial a_2}{\partial\vh_1}
    \cdot \frac{\partial\vh_1}{\partial a_1}
    \cdot \frac{\partial a_1}{\partial W}
\]
The bracketed terms (in boxes in the lecture notes) are \textbf{reused} from the decoder backprop pass; the remaining terms are \textbf{recomputed} for the encoder.

%=====================================================================
\section{Link Between PCA and Autoencoders}
%=====================================================================

\subsection{Setup: When Are They Equivalent?}

\begin{theorembox}
The encoder of a linear autoencoder is \textbf{equivalent to PCA} when:
\begin{enumerate}
  \item The encoder is \textbf{linear}: $\vh = W\vx$ (not $\sigma(W\vx)$ or $\tanh(W\vx)$)
  \item The decoder is \textbf{linear}: $\hatx = W^*\vh + \vc$
  \item The loss function is \textbf{mean squared error}
  \item Inputs are \textbf{normalised} to zero mean by:
    \[
      \hat{x}_{ij} = \frac{1}{\sqrt{m}}\left(x_{ij} - \frac{1}{m}\sum_{k=1}^{m}x_{kj}\right)
    \]
\end{enumerate}
\end{theorembox}

\paragraph{Why normalise?} After normalisation, $X^\top X = \frac{1}{m}(X')^\top X'$ is the \textbf{covariance matrix}, which plays a central role in PCA.

%-------------------------------------------------------------------
\subsection{The Optimisation Problem}
%-------------------------------------------------------------------

With linear encoder and decoder (ignoring biases), the autoencoder objective becomes:
\[
  \min_{W,W^*} \norm{X - HW^*}_F^2
\]
where:
\begin{itemize}
  \item $X \in \R^{m\times n}$: data matrix (rows = samples)
  \item $H = XW^\top \in \R^{m\times k}$: hidden representation matrix
  \item $W \in \R^{k\times n}$: encoder weights
  \item $W^* \in \R^{k\times n}$: decoder weights
  \item $\hat{X} = HW^* \in \R^{m\times n}$: reconstructed data
  \item $\norm{A}_F = \sqrt{\sum_{ij}a_{ij}^2}$ is the \textbf{Frobenius norm}
\end{itemize}

%-------------------------------------------------------------------
\subsection{Connection to SVD (Eckart--Young Theorem)}
%-------------------------------------------------------------------

Writing $\hat{X} = HW^* = U\Sigma V^\top$ in matrix form and using the Frobenius norm definition, the \textbf{Eckart--Young--Mirsky theorem} from SVD tells us:

\begin{theorembox}[title={Eckart--Young Theorem}]
The best rank-$k$ approximation to $X$ in Frobenius norm is:
\[
  \hat{X} = HW^* = U_{\cdot,\leq k}\,\Sigma_{k\times k}\,V_{\cdot,\leq k}^\top
\]
where $X = U\Sigma V^\top$ is the full SVD.
\end{theorembox}

%-------------------------------------------------------------------
\subsection{Matching Variables: Finding $W$ and $W^*$}
%-------------------------------------------------------------------

By matching variables in $HW^* = U_{\cdot,\leq k}\Sigma_{k\times k}V_{\cdot,\leq k}^\top$, one possible solution is:
\[
  H = U_{\cdot,\leq k}\Sigma_{k\times k},
  \qquad
  W^* = V_{\cdot,\leq k}^\top
\]
Other solutions also exist (any rotation in the latent space), e.g.:
\[
  H = U_{\cdot,\leq k},
  \qquad
  W^* = \Sigma_{k\times k}V_{\cdot,\leq k}^\top
\]

%-------------------------------------------------------------------
\subsection{Full Derivation of the Encoder $W$}
%-------------------------------------------------------------------

Working with the first solution, we want to show $H = WX$ (linear encoding). We have $H = U_{\cdot,\leq k}\Sigma_{k\times k}$ and $X = U\Sigma V^\top$.

\textbf{Step 1:} Pre-multiply by $(XX^\top)(XX^\top)^{-1} = I$:
\begin{align*}
  H &= (XX^\top)(XX^\top)^{-1} \cdot U_{\cdot,\leq k}\Sigma_{k\times k}
\end{align*}

\textbf{Step 2:} Compute $XX^\top$ using $X = U\Sigma V^\top$:
\[
  XX^\top = U\Sigma V^\top \cdot V\Sigma^\top U^\top = U\Sigma\Sigma^\top U^\top = U\Sigma^2 U^\top
\]
(since $V^\top V = I$ and $\Sigma$ is rectangular diagonal)

\textbf{Step 3:} Substitute into $H$ expression:
\begin{align*}
  H &= (U\Sigma^2 U^\top)(U\Sigma^2 U^\top)^{-1} U_{\cdot,\leq k}\Sigma_{k\times k}
\end{align*}

\textbf{Step 4:} Expand using $X = U\Sigma V^\top$, so $X^\top = V\Sigma^\top U^\top$:
\begin{align*}
  H &= XV\Sigma^\top U^\top (U\Sigma\Sigma^\top U^\top)^{-1} U_{\cdot,\leq k}\Sigma_{k\times k}\\
    &= XV\Sigma^\top (U^\top U)(\Sigma\Sigma^\top)^{-1} U^\top U_{\cdot,\leq k}\Sigma_{k\times k}
\end{align*}

Since $U^\top U = I$ and $U^\top U_{\cdot,\leq k} = I_{\cdot,\leq k}$ (only first $k$ columns survive):
\begin{align*}
  H &= XV\underbrace{\Sigma^\top(\Sigma\Sigma^\top)^{-1}}_{\Sigma^{-1}} I_{\cdot,\leq k}\Sigma_{k\times k}\\
    &= XV\Sigma^{-1}I_{\cdot,\leq k}\Sigma_{k\times k}\\
    &= XV\,I_{\cdot,\leq k}\\
    &= XV_{\cdot,\leq k}
\end{align*}

\textbf{Step 5:} So $H = XV_{\cdot,\leq k}$, which means $W = V_{\cdot,\leq k}^\top$ (since $H = WX$ requires $W = V_{\cdot,\leq k}^\top$).

\begin{theorembox}
\[
  \boxed{H = XV_{\cdot,\leq k},\quad W = V_{\cdot,\leq k}^\top}
\]
$H$ is a linear transformation of $X$ -- hence the optimal encoder is a linear encoder.
\end{theorembox}

%-------------------------------------------------------------------
\subsection{Conclusion: PCA--AE Equivalence}
%-------------------------------------------------------------------

The optimal linear autoencoder gives:
\begin{itemize}
  \item \textbf{Encoder}: $W = V_{\cdot,\leq k}^\top$ (first $k$ right singular vectors of $X$)
  \item \textbf{Hidden representation}: $H = XV_{\cdot,\leq k}$ (projection of data onto principal components)
  \item \textbf{Decoder}: $W^* = V_{\cdot,\leq k}^\top$
\end{itemize}

\begin{theorembox}[title={PCA--AE Equivalence Proof}]
From SVD: $V$ is the matrix of eigenvectors of $X^\top X$.\\
From PCA: $P$ is the matrix of eigenvectors of the covariance matrix.\\
After normalisation: $X^\top X = \frac{1}{m}(X')^\top X'$ is exactly the covariance matrix.\\
$\Rightarrow$ Encoder matrix $W$ for linear AE $=$ projection matrix $P$ for PCA. \hfill\textbf{QED}
\end{theorembox}

\begin{notebox}
\textbf{Summary: Remember These Four Conditions}
\begin{enumerate}
  \item Use a \textbf{linear encoder} ($g,f$ are linear functions)
  \item Use a \textbf{linear decoder}
  \item Use \textbf{squared-error loss}
  \item \textbf{Normalise the inputs} to zero mean: $\hat{x}_{ij} = \frac{1}{\sqrt{m}}\left(x_{ij} - \frac{1}{m}\sum_{k=1}^{m}x_{kj}\right)$
\end{enumerate}
\end{notebox}

%=====================================================================
\section{Regularisation on Autoencoders}
%=====================================================================

\subsection{Why Regularisation Is Needed}

\begin{alertbox}
\textbf{Overfitting in Autoencoders:}
\begin{itemize}
  \item More learnable parameters $\Rightarrow$ more overfitting
  \item Over-complete AEs: the model can simply copy $\vx_i$ into $\vh$ and $\vh$ into $\hatx_i$ (identity function -- useless)
  \item Even under-complete AEs: with very high-dimensional inputs (e.g.\ $100\times100$ images, $k=1000$), there are still many parameters and much redundancy in the input $\Rightarrow$ overfitting
\end{itemize}
Any kind of regularisation always tries to \textbf{restrict the freedom of parameters}.
\end{alertbox}

%-------------------------------------------------------------------
\subsection{L2 Regularisation (Weight Decay)}
%-------------------------------------------------------------------

The simplest solution: add an L2 penalty term to the objective:
\[
  \min_{\theta}\;
  \underbrace{\frac{1}{m}\sum_{i=1}^{m}\sum_{j=1}^{n}(\hat{x}_{ij} - x_{ij})^2}_{\text{reconstruction loss}}
  + \underbrace{\lambda\norm{\theta}^2}_{\text{L2 penalty}}
\]

\subsubsection{Gradient with L2 Regularisation}

For $\tilde{\Loss}(w) = \Loss(w) + \frac{\alpha}{2}\norm{w}^2$:
\[
  \grad\tilde{\Loss}(w) = \grad\Loss(w) + \alpha w
\]

The SGD update rule becomes:
\[
  w_{t+1} = w_t - \eta\grad\tilde{\Loss}(w_t)
  = w_t - \eta\grad\Loss(w_t) - \eta\alpha w_t
  = (1 - \eta\alpha)w_t - \eta\grad\Loss(w_t)
\]

This is very easy to implement: just add $\alpha w$ to every gradient $\frac{\partial\Loss}{\partial w}$.

\subsubsection{Geometric Interpretation via Taylor Expansion}

Let $w^*$ be the unregularised optimal solution: $\grad\Loss(w^*) = 0$.

Consider $w = w^* + h$ in the neighbourhood of $w^*$. Using Taylor series to second order:
\[
  \Loss(w^* + h) \approx \Loss(w^*)
    + h^\top \underbrace{\grad\Loss(w^*)}_{=0}
    + \tfrac{1}{2}h^\top H h
\]
\[
  \Loss(w) \approx \Loss(w^*) + \tfrac{1}{2}(w-w^*)^\top H (w-w^*)
\]

where $H$ is the Hessian. Therefore:
\[
  \boxed{\grad\Loss(w) \approx H(w - w^*)}
\]

\subsubsection{Finding the Regularised Solution}

The regularised loss gradient:
\[
  \grad\tilde{\Loss}(w) = \grad\Loss(w) + \alpha w \approx H(w-w^*) + \alpha w
\]

Setting $\grad\tilde{\Loss}(\tilde{w}) = 0$ at the regularised optimum $\tilde{w}$:
\begin{align*}
  H(\tilde{w} - w^*) + \alpha\tilde{w} &= 0\\
  H\tilde{w} - Hw^* + \alpha\tilde{w} &= 0\\
  (H + \alpha I)\tilde{w} &= Hw^*\\
  \tilde{w} &= (H + \alpha I)^{-1}Hw^*
\end{align*}

\textbf{Note:} As $\alpha \to 0$, $\tilde{w} \to w^*$ (recovers unregularised solution).

\subsubsection{Spectral Analysis of L2 Regularisation}

Since $H$ is symmetric PSD, it has spectral decomposition $H = Q\Lambda Q^\top$ where $Q$ is orthogonal ($QQ^\top = Q^\top Q = I$) and $\Lambda = \mathrm{diag}(\lambda_1, \ldots, \lambda_n)$.

\begin{align*}
  \tilde{w}
  &= (Q\Lambda Q^\top + \alpha I)^{-1}Q\Lambda Q^\top w^*\\
  &= (Q\Lambda Q^\top + \alpha QIQ^\top)^{-1}Q\Lambda Q^\top w^*\\
  &= \bigl(Q(\Lambda + \alpha I)Q^\top\bigr)^{-1}Q\Lambda Q^\top w^*\\
  &= (Q^\top)^{-1}(\Lambda + \alpha I)^{-1}Q^{-1} \cdot Q\Lambda Q^\top w^*\\
  &= Q(\Lambda + \alpha I)^{-1}\Lambda Q^\top w^*
\end{align*}

\[
  \boxed{\tilde{w} = Q\,D\,Q^\top w^*}
\]

where $D = (\Lambda + \alpha I)^{-1}\Lambda = \mathrm{diag}\!\left(\frac{\lambda_1}{\lambda_1+\alpha},\ldots,\frac{\lambda_n}{\lambda_n+\alpha}\right)$.

\subsubsection{What Does $D$ Mean?}

\[
  D_{ii} = \frac{\lambda_i}{\lambda_i + \alpha}
\]

The operation $\tilde{w} = QDQ^\top w^*$ has a clear geometric interpretation:
\begin{enumerate}
  \item $w^*$ is first \textbf{rotated} by $Q^\top$ to give $Q^\top w^*$
  \item Each element $i$ is \textbf{scaled} by $\frac{\lambda_i}{\lambda_i + \alpha}$
  \item The result is \textbf{rotated back} by $Q$
\end{enumerate}

\begin{itemize}
  \item If $\lambda_i \gg \alpha$: $\frac{\lambda_i}{\lambda_i+\alpha} \approx 1$ -- direction almost unchanged (important direction retained)
  \item If $\lambda_i \ll \alpha$: $\frac{\lambda_i}{\lambda_i+\alpha} \approx 0$ -- direction shrunk to zero (unimportant direction removed)
\end{itemize}

\begin{keybox}
\textbf{Only significant directions (larger eigenvalues) are retained.} Less important weights shrink, more important weights stay.

\textbf{Effective number of parameters} $= \displaystyle\sum_{i=1}^{n}\frac{\lambda_i}{\lambda_i + \alpha}$
\end{keybox}

%-------------------------------------------------------------------
\subsection{Weight Tying}
%-------------------------------------------------------------------

\begin{definitionbox}
\textbf{Weight tying} enforces $W^* = W^\top$, i.e.\ the decoder weight matrix is the transpose of the encoder weight matrix.
\end{definitionbox}

By tying, we reduce the number of free parameters from $2d\times n$ to $d\times n$ -- exactly half. This acts as regularisation by reducing the model's freedom.

\paragraph{Why $W^\top$ instead of $W^{-1}$?}
\begin{itemize}
  \item The encoder projects onto feature directions
  \item The decoder reconstructs from the \emph{same} directions (coherent)
  \item Using $W^{-1}$ would require $W$ to be square and invertible -- a strong and unnecessary constraint
\end{itemize}

\begin{alertbox}
\textbf{Weight Tying Summary:} Tie weight learning so that the same weight works at both encoder and decoder. This is commonly used for regularisation. It imposes an equality constraint between parameters, halving the total number of free parameters.
\end{alertbox}

%=====================================================================
\section{Denoising Autoencoders}
%=====================================================================

\subsection{Motivation and Core Idea}

A \textbf{denoising autoencoder} (DAE) corrupts the input before feeding it to the network, but still tries to reconstruct the \emph{original} uncorrupted input.

\paragraph{Corruption process:} A probabilistic process $P(\tilde{x}_{ij} \mid x_{ij})$ generates a corrupted version $\tilde{\vx}_i$ of the original input $\vx_i$.

\paragraph{Objective:} Still minimise reconstruction of the \emph{original uncorrupted} input:
\[
  \arg\min_{\theta}\;\frac{1}{m}\sum_{i=1}^{m}\sum_{j=1}^{n}(\hat{x}_{ij} - x_{ij})^2
\]
but feeding $\tilde{\vx}_i$ (corrupted) to the network, not $\vx_i$ (original).

\begin{keybox}
\textbf{Why does corruption help?} The network cannot simply copy input to output (it receives $\tilde{\vx}_i$ but must output $\vx_i$). It \emph{must} learn the true structure of the data. This is a ``more robust'' kind of regularisation.
\end{keybox}

%-------------------------------------------------------------------
\subsection{Corruption Processes}
%-------------------------------------------------------------------

\subsubsection{Masking Noise}
Set each input to 0 with probability $q$:
\[
  P(\tilde{x}_{ij} = 0 \mid x_{ij}) = q,
  \quad
  P(\tilde{x}_{ij} = x_{ij} \mid x_{ij}) = 1 - q
\]
With probability $q$, the input is ``masked''; with probability $1-q$, it is retained as-is.

\subsubsection{Gaussian Noise}
Add Gaussian noise:
\[
  \tilde{x}_{ij} = x_{ij} + \mathcal{N}(0, \sigma^2)
\]

\begin{alertbox}
With corruption, it no longer makes sense for the model to copy $\tilde{\vx}_i$ into $\vh$ and then into $\hatx_i$ -- the objective function will not be minimised by doing so. Instead, the model must capture the actual characteristics of the data correctly.

For example, to reconstruct a corrupted feature $\tilde{x}_{ij}$, the network must rely on its \emph{interactions} with other features of $\vx_i$ -- learning true correlations in the data.
\end{alertbox}

%-------------------------------------------------------------------
\subsection{Visualising Hidden Neurons in Denoising AEs}
%-------------------------------------------------------------------

We can think of each hidden neuron as a \textbf{filter} which maximally activates for certain input configurations.

\subsubsection{Finding the Maximally Activating Input}

For neuron $\ell$ with $h_\ell = \sigma(w_\ell^\top\vx_i)$, after training, find what input $\vx_i$ causes $h_\ell$ to be maximum:
\[
  \max_{\vx_i}\; w_\ell^\top\vx_i
  \quad\text{subject to}\quad
  \norm{\vx_i}^2 = 1
\]

By the Cauchy-Schwarz inequality: $w^\top x \leq \norm{w}\norm{x}$, with equality when $x = cw$ for some scalar $c > 0$.

From the constraint $\norm{x}^2 = 1$:
\[
  1 = \norm{cw}^2 = c^2\norm{w}^2
  \quad\Rightarrow\quad
  c = \frac{1}{\norm{w}}
\]

So the maximally activating input is:
\[
  \boxed{\vx_i = \frac{w_\ell}{\sqrt{w_\ell^\top w_\ell}} = \frac{w_\ell}{\norm{w_\ell}}}
\]

The \textbf{normalised weight vector} of each neuron corresponds to its maximally activating input.

\subsubsection{Visual Comparison: Vanilla AE vs.\ Denoising AE}

We can plot the images $\vx_i = w_\ell/\norm{w_\ell}$ (using weights $w_1, w_2, \ldots, w_k$ learned by each AE) to see what each neuron ``looks for'':

\begin{itemize}
  \item \textbf{Vanilla AE} (no noise): Does not learn many meaningful patterns; neurons look noisy or random
  \item \textbf{Denoising AE with $q = 0.25$}: Learns more structured patterns; neurons resemble stroke-detectors
  \item \textbf{Denoising AE with $q = 0.5$}: Thicker filters. As noise increases, the neuron must rely on \emph{adjacent pixels} to feel confident about a stroke, so filters become wider (more neighbourhood information)
\end{itemize}

\begin{keybox}
Hidden neurons of the denoising AE learn to act like \textbf{pen-stroke detectors}. As the noise level $q$ increases:
\begin{itemize}
  \item Filters become wider (larger receptive field)
  \item Neurons become more robust and neighbourhood-aware
  \item Weights become larger -- everything is more generalised
\end{itemize}
\end{keybox}

%-------------------------------------------------------------------
\subsection{Application: Handwritten Digit Recognition (MNIST)}
%-------------------------------------------------------------------

Input: $\vx_i \in \R^{784}$ ($28\times28$ pixels).

\paragraph{Basic approach:} Train directly with 9 output neurons (digits 0--9).

\paragraph{AE approach:}
\begin{enumerate}
  \item First learn important characteristics of data using AE: $\vx_i \in \R^{784} \to \vh \in \R^9$
  \item Produce better and better representations through stacked AEs
  \item Feed learned features $\vh$ to a multi-class SVM
\end{enumerate}

\paragraph{Advantages of AE approach:}
\begin{itemize}
  \item Better feature representation than raw pixels
  \item Smaller representation ($\vh \in \R^d$ instead of $\R^{784}$)
  \item No feature engineering needed -- AE learns features automatically
  \item Can use SVM more effectively on learned representations
\end{itemize}

\paragraph{Disadvantage:} \textbf{No explainability} -- the learned features are not directly interpretable.

%=====================================================================
\section{Sparse Autoencoders}
%=====================================================================

\subsection{Motivation: Sparsity Constraint}

A \textbf{sparse autoencoder} (SAE) tries to ensure each hidden neuron is \textbf{inactive most of the time}.

\begin{definitionbox}
With sigmoid activation, hidden neurons have values in $[0,1]$:
\begin{itemize}
  \item \textbf{Activated}: output close to 1
  \item \textbf{Not activated}: output close to 0 (sparse)
\end{itemize}

The \textbf{average activation} of neuron $\ell$ over the dataset:
\[
  \hat{\rho}_\ell = \frac{1}{m}\sum_{i=1}^{m}h(\vx_i)_\ell
\]

A sparse autoencoder uses a \textbf{sparsity parameter} $\rho \approx 0.005$ and enforces:
\[
  \boxed{\hat{\rho}_\ell = \rho}
\]
i.e.\ the average activation should equal $\rho$ (near zero).
\end{definitionbox}

\begin{keybox}
\textbf{Intuition:} Overfitting happens because of too much freedom. Sparsity restricts this freedom. If a neuron fires for everything (e.g.\ $\hat{\rho}_\ell \approx 1$), there's no selective power. By forcing it to fire only for some patterns, we build more meaningful representations.
\end{keybox}

%-------------------------------------------------------------------
\subsection{KL Divergence Penalty}
%-------------------------------------------------------------------

To enforce $\hat{\rho}_\ell = \rho$, add the following regularisation term (sum of KL divergences):
\[
  \Reg(\theta)
  = \sum_{\ell=1}^{k}
    \left[
      \rho\log\frac{\rho}{\hat{\rho}_\ell}
      + (1-\rho)\log\frac{1-\rho}{1-\hat{\rho}_\ell}
    \right]
\]

The full loss becomes:
\[
  \tilde{\Loss}(\theta) = \Loss(\theta) + \Reg(\theta)
\]
where $\Loss(\theta)$ is the squared-error or cross-entropy reconstruction loss.

\subsubsection{When Does $\Reg(\theta)$ Reach Its Minimum?}

$\Reg(\theta)$ is a sum of \textbf{KL divergences} between a Bernoulli$(\rho)$ and a Bernoulli$(\hat{\rho}_\ell)$:
\[
  \text{KL}(\text{Bern}(\rho)\;\|\;\text{Bern}(\hat{\rho}_\ell))
  = \rho\log\frac{\rho}{\hat{\rho}_\ell} + (1-\rho)\log\frac{1-\rho}{1-\hat{\rho}_\ell}
\]

KL divergence is always $\geq 0$, with equality iff the two distributions are identical.

\[
  \Reg(\theta) = 0 \iff \hat{\rho}_\ell = \rho \;\text{ for all } \ell
\]

So $\Reg(\theta)$ is minimised exactly when the sparsity constraint is satisfied.

\begin{notebox}
\textbf{KL Divergence Intuition:} ``How much information is lost when you use one probability distribution to approximate another.'' Only a few neurons are allowed to be active at any time. Even with many hidden neurons, sparsity prevents the network from simply memorising.
\end{notebox}

%-------------------------------------------------------------------
\subsection{Gradient of the Sparsity Term}
%-------------------------------------------------------------------

To use backpropagation, we need $\frac{\partial\Reg(\theta)}{\partial w}$.

First, derivative of $\Reg$ with respect to $\hat{\rho}_\ell$:
\[
  \frac{\partial\Reg}{\partial\hat{\rho}_\ell}
  = -\frac{\rho}{\hat{\rho}_\ell} + \frac{1-\rho}{1-\hat{\rho}_\ell}
\]

For a single element $w_{j\ell}$ of the weight matrix:
\begin{align*}
  \frac{\partial\hat{\rho}_\ell}{\partial w_{j\ell}}
  &= \frac{\partial}{\partial w_{j\ell}}\left[\frac{1}{m}\sum_{i=1}^{m}g(w_j^\top\vx_i + b_j)\right]
  = \frac{1}{m}\sum_{i=1}^{m}g'(w_j^\top\vx_i + b_j)\cdot x_{ij}
\end{align*}

In matrix notation:
\[
  \frac{\partial\hat{\rho}_\ell}{\partial w}
  = \vx_i\bigl(g'(w^\top\vx_i + \vb)\bigr)^\top
\]

By chain rule:
\[
  \frac{\partial\Reg(\theta)}{\partial w}
  = \frac{\partial\Reg(\theta)}{\partial\hat{\rho}_\ell}
    \cdot \frac{\partial\hat{\rho}_\ell}{\partial w}
  = \left(-\frac{\rho}{\hat{\rho}_\ell} + \frac{1-\rho}{1-\hat{\rho}_\ell}\right)
    \cdot \vx_i\bigl(g'(w^\top\vx_i + \vb)\bigr)^\top
\]

Final update (adding both terms):
\[
  \frac{\partial\tilde{\Loss}(\theta)}{\partial w}
  = \frac{\partial\Loss(\theta)}{\partial w}
  + \frac{\partial\Reg(\theta)}{\partial w}
\]

%=====================================================================
\section{Contractive Autoencoders}
%=====================================================================

\subsection{Motivation}

A \textbf{contractive autoencoder} (CAE) also prevents over-complete AEs from learning the identity function, but in a completely different way.

It penalises the \textbf{sensitivity} of the hidden representation to changes in the input -- i.e.\ it adds the Frobenius norm of the Jacobian:

\[
  \Reg(\theta) = \norm{J_f(\vh)}_F^2
  = \sum_{j=1}^{n}\sum_{\ell=1}^{k}
    \left(\frac{\partial h_\ell}{\partial x_j}\right)^2
\]

where $\vh$ is the hidden representation (size $k$) and $J_f(\vh) \in \R^{k\times n}$ is the Jacobian of $\vh$ with respect to $\vx$.

%-------------------------------------------------------------------
\subsection{The Jacobian Matrix}
%-------------------------------------------------------------------

\[
  J_f(\vh)
  = \begin{bmatrix}
      \frac{\partial h_1}{\partial x_1} & \cdots & \frac{\partial h_1}{\partial x_n}\\
      \frac{\partial h_2}{\partial x_1} & \cdots & \frac{\partial h_2}{\partial x_n}\\
      \vdots & \ddots & \vdots\\
      \frac{\partial h_k}{\partial x_1} & \cdots & \frac{\partial h_k}{\partial x_n}
    \end{bmatrix}
  \in \R^{k\times n}
\]

The $(i,j)$ entry captures the variation in the output of the $i$-th neuron with a small variation in the $j$-th input. It measures \textbf{sensitivity}.

%-------------------------------------------------------------------
\subsection{Intuition: A Beautiful Contradiction}
%-------------------------------------------------------------------

Consider $\frac{\partial h_1}{\partial x_1} = 0$: neuron $h_1$ is \emph{not sensitive} to changes in $x_1$.

\begin{itemize}
  \item \textbf{From $\Loss(\theta)$ perspective}: $\vh$ must capture variations in $\vx$ to reconstruct it. So we \emph{want} sensitivity.
  \item \textbf{From $\Reg(\theta)$ perspective}: We want $\vh$ to be \emph{insensitive} to $\vx$.
\end{itemize}

These are \textbf{contradicting objectives}, and that is exactly the point!

\begin{keybox}
\textbf{The Contradiction That Creates Good Representations:}\\
By placing two contradicting objectives against each other, we ensure $\vh$ is sensitive to \emph{only the most important variations} as observed in the training data.

\begin{itemize}
  \item $\Loss(\theta)$: captures \textbf{important} variations in data
  \item $\Reg(\theta)$ (Jacobian penalty): penalises sensitivity to \emph{any} variation
  \item \textbf{Trade-off}: capture only the \emph{most important} variations
\end{itemize}
\end{keybox}

%-------------------------------------------------------------------
\subsection{Analogy with PCA}
%-------------------------------------------------------------------

Consider data varying along directions $u_1$ (high variance) and $u_2$ (low variance/noise).
\begin{itemize}
  \item It makes sense to maximise a neuron's sensitivity to variations along $u_1$ (large variance direction)
  \item It makes sense to inhibit sensitivity along $u_2$ (small variations -- likely noise, unimportant for reconstruction)
\end{itemize}

This is exactly what PCA does: keep the high-variance directions and discard the rest.

%-------------------------------------------------------------------
\subsection{Frobenius Norm for Sigmoid Units}
%-------------------------------------------------------------------

For sigmoid hidden units $h_\ell = \sigma(w_\ell^\top\vx + b_\ell)$:
\[
  \frac{\partial h_\ell}{\partial x_j}
  = h_\ell(1-h_\ell)\cdot w_{\ell j}
\]

Therefore:
\[
  \norm{J_f(\vh)}_F^2
  = \sum_\ell (h_\ell(1-h_\ell))^2 \norm{w_\ell}^2
\]

The combined objective:
\[
  \min_{\theta}\;\Loss(\theta) + \lambda\norm{J_f(\vh)}_F^2
  = \min_{\theta}\;\Loss(\theta) + \lambda\sum_\ell(h_\ell(1-h_\ell))^2\norm{w_\ell}^2
\]

%-------------------------------------------------------------------
\subsection{Regularisation Methods Comparison}
%-------------------------------------------------------------------

\begin{center}
\begin{tabular}{lll}
\toprule
\textbf{Method} & \textbf{Regularisation Term $\Reg(\theta)$} & \textbf{What It Controls}\\
\midrule
L2 / Weight Decay & $\lambda\norm{\theta}^2$ & Weight magnitude \\
Sparse AE (KL) & $\sum_\ell\mathrm{KL}(\rho\|\hat{\rho}_\ell)$ & Average activation \\
Contractive AE & $\norm{J_f(\vh)}_F^2$ & Sensitivity to input \\
Denoising AE & Corruption $P(\tilde{x}\mid x)$ & Implicit (via corruption) \\
Weight Tying & $W^* = W^\top$ & Parameter count \\
\bottomrule
\end{tabular}
\end{center}

%=====================================================================
\section{Bias, Variance and Model Complexity}
%=====================================================================

\subsection{Experimental Setup: Fitting a Sinusoidal Function}

Consider fitting a sinusoidal function to a set of 25 data points with two models:
\begin{itemize}
  \item \textbf{Simple model}: $y = \hat{f}(x) = w_1 x + w_0$ (degree 1 polynomial)
  \item \textbf{Complex model}: $y = \hat{f}(x) = \sum_{i=1}^{25}w_i x^i + w_0$ (degree 25 polynomial)
\end{itemize}

We repeat the experiment $K$ times with different samples of 25 training points.

\paragraph{Observations:}
\begin{itemize}
  \item \textbf{Simple models} trained on different samples: Do not differ much from each other. However, they are very far from the true sinusoidal curve. $\Rightarrow$ \textbf{High bias, low variance (underfitting)}
  \item \textbf{Complex models} trained on different samples: Are very different from each other (each overfits to its training data differently). $\Rightarrow$ \textbf{Low bias, high variance (overfitting)}
\end{itemize}

%-------------------------------------------------------------------
\subsection{Formal Definitions}
%-------------------------------------------------------------------

Let $f(x)$ be the true model and $\hat{f}(x)$ be our estimate trained on a particular dataset.

\begin{definitionbox}
The \textbf{bias} of $\hat{f}(x)$ measures systematic error:
\[
  \Bias(\hat{f}(x)) = \E[\hat{f}(x)] - f(x)
\]
where $\E[\hat{f}(x)]$ is the expected (average) value of the model over all possible training sets.

The \textbf{variance} of $\hat{f}(x)$ measures spread/sensitivity:
\[
  \Var(\hat{f}(x)) = \E\!\left[\bigl(\hat{f}(x) - \E[\hat{f}(x)]\bigr)^2\right]
\]
It tells us how much the different $\hat{f}(x)$'s trained on different samples spread out.
\end{definitionbox}

\begin{itemize}
  \item Simple model: high bias (average blue line far from $f(x)$), low variance (lines don't differ much)
  \item Complex model: low bias (average close to $f(x)$), high variance (lines differ greatly)
\end{itemize}

Informally: simple model = high bias, low variance; complex model = low bias, high variance. There is always a \textbf{trade-off}.

%-------------------------------------------------------------------
\subsection{Bias--Variance Decomposition: Full Derivation}
%-------------------------------------------------------------------

\begin{theorembox}[title={Bias--Variance Decomposition}]
For any data point $(x, y)$ with $y = f(x) + \varepsilon$ where $\varepsilon \sim \mathcal{N}(0, \sigma^2)$:
\[
  \E\!\left[(y - \hat{f}(x))^2\right]
  = \bigl(\Bias(\hat{f}(x))\bigr)^2 + \Var(\hat{f}(x)) + \sigma^2
\]
\end{theorembox}

\paragraph{Proof:} Define:
\[
  A = f(x) - \E[\hat{f}(x)],\quad
  B = \varepsilon,\quad
  C = \E[\hat{f}(x)] - \hat{f}(x)
\]

Note: $y - \hat{f}(x) = f(x) + \varepsilon - \hat{f}(x) = A + B + C$

\begin{align*}
  \E[(y-\hat{f})^2]
  &= \E\!\left[(A+B+C)^2\right]\\
  &= \E[A^2] + \E[B^2] + \E[C^2]
    + 2\E[AB] + 2\E[AC] + 2\E[BC]
\end{align*}

\textbf{Cross term $\E[AB]$:}
\[
  \E[AB] = \E[(f - \E[\hat{f}])\varepsilon]
  = (f-\E[\hat{f}])\E[\varepsilon] = 0
\]
since $\varepsilon$ has zero mean and is independent of $f - \E[\hat{f}]$ (which is a constant).

\textbf{Cross term $\E[BC]$:}
\[
  \E[BC] = \E[\varepsilon(\E[\hat{f}] - \hat{f})]
  = \E[\varepsilon]\cdot\E[\E[\hat{f}]-\hat{f}] = 0
\]
since $\E[\varepsilon]=0$ and $\varepsilon$ is independent of the model randomness.

\textbf{Cross term $\E[AC]$:}
\[
  \E[AC] = \E[(f-\E[\hat{f}])(\E[\hat{f}]-\hat{f})]
\]
Since $f - \E[\hat{f}]$ is a constant (with respect to model randomness -- it does not depend on the particular training set), and $\E[\E[\hat{f}]-\hat{f}] = \E[\hat{f}] - \E[\hat{f}] = 0$:
\[
  \E[AC] = (f - \E[\hat{f}])\cdot\underbrace{\E[\E[\hat{f}]-\hat{f}]}_{=0} = 0
\]

\textbf{Remaining terms:}
\begin{align*}
  \E[A^2] &= (f - \E[\hat{f}])^2 = \bigl(\Bias(\hat{f})\bigr)^2\\
  \E[B^2] &= \E[\varepsilon^2] = \Var(\varepsilon) = \sigma^2\\
  \E[C^2] &= \E\!\left[(\E[\hat{f}] - \hat{f})^2\right] = \Var(\hat{f})
\end{align*}

Therefore:
\[
  \boxed{
    \E\!\left[(y - \hat{f}(x))^2\right]
    = \underbrace{\bigl(\Bias(\hat{f}(x))\bigr)^2}_{\substack{\text{systematic error}\\\text{wrong assumptions}}}
    + \underbrace{\Var(\hat{f}(x))}_{\substack{\text{error from}\\\text{sensitivity to training}}}
    + \underbrace{\sigma^2}_{\substack{\text{irreducible}\\\text{noise}}}
  }
\]

\hfill\textbf{QED}

%=====================================================================
\section{Training Error vs.\ Test Error}
%=====================================================================

\subsection{Setup}

The parameters of $\hat{f}(x)$ are trained using training set $\{(x_i, y_i)\}_{i=1}^{n}$. At test time, we evaluate on an \emph{unseen} validation set.

\begin{definitionbox}
\begin{itemize}
  \item \textbf{Training error}: $\frac{1}{n}\sum_{i=1}^{n}(y_i - \hat{f}(x_i))^2$
  \item \textbf{Test error} (MSE on new points): $\frac{1}{n_{\text{test}}}\sum_{i=n+1}^{n+n_{\text{test}}}(y_i - \hat{f}(x_i))^2$
\end{itemize}
\end{definitionbox}

As model complexity increases, training error becomes overly optimistic. The validation error gives the real picture.

\textbf{Bias-Variance U-curve:}
\begin{itemize}
  \item High bias region (left): model too simple, high train and test error
  \item Sweet spot (middle): best model complexity -- minimises test error
  \item High variance region (right): model too complex, low train error but high test error
\end{itemize}

%-------------------------------------------------------------------
\subsection{Why Test Error Approximates True Error}
%-------------------------------------------------------------------

We want to estimate (true expectation):
\[
  \E\!\left[(\hat{f}(x_i) - f(x_i))^2\right]
\]

But we cannot compute this directly since we don't know $f$. Let $y_i = f(x_i) + \varepsilon_i$.

For any data point: $\hat{y}_i - y_i = \hat{f}(x_i) - f(x_i) - \varepsilon_i$

\begin{align*}
  \E[(\hat{f}(x_i) - f(x_i))^2]
  &= \E[(\hat{y}_i - y_i)^2] - \E[\varepsilon_i^2]
    + 2\E[\varepsilon_i(\hat{f}(x_i) - f(x_i))]
\end{align*}

\subsubsection{Case 1: Using Test Observations}

Test observations were \emph{not} used to estimate parameters of $\hat{f}$. Hence $\varepsilon_i \perp \hat{f}(x_i)$:
\[
  \E[\varepsilon_i(\hat{f}(x_i) - f(x_i))]
  = \E[\varepsilon_i]\cdot\E[\hat{f}(x_i) - f(x_i)]
  = 0 \cdot \E[\hat{f}(x_i) - f(x_i)]
  = 0
\]

Therefore:
\[
  \boxed{
    \E[(\hat{f}(x_i) - f(x_i))^2]
    \approx \text{Empirical test error} + \text{small constant}
  }
\]

The empirical test error is a \textbf{true picture} of prediction error (not overly optimistic).

\subsubsection{Case 2: Using Training Observations}

Training observations \emph{were} used to estimate parameters. Hence $\varepsilon_i$ is \emph{not} independent of $\hat{f}(x_i)$:
\[
  \E[\varepsilon_i(\hat{f}(x_i) - f(x_i))]
  \neq \E[\varepsilon_i]\cdot\E[\hat{f}(x_i) - f(x_i)]
\]

The extra covariance term is \emph{not} zero:
\[
  \Cov(\varepsilon_i,\; \hat{f}(x_i) - f(x_i)) \neq 0
\]

So $\E[(\hat{f}(x_i)-f(x_i))^2] \neq \text{empirical training error}$.

The empirical training error is \textbf{smaller} than the true error -- the model is \textbf{optimistic}!

\begin{alertbox}
\textbf{Always use a validation set (independent of training set) to estimate error.} This gives the closest approximation to true error. Training error is biased downward.
\end{alertbox}

%=====================================================================
\section{True Error and Model Complexity: Stein's Lemma}
%=====================================================================

\subsection{Stein's Lemma}

\begin{theorembox}[title={Stein's Lemma}]
If $X \sim \mathcal{N}(\mu, \sigma^2)$ and $g$ is a differentiable function, then:
\[
  \E[(X-\mu)g(X)] = \sigma^2\,\E[g'(X)]
\]
Equivalently, for $X \sim \mathcal{N}(0, \sigma^2)$:
\[
  \E[Xg(X)] = \sigma^2\,\E[g'(X)]
\]
\end{theorembox}

%-------------------------------------------------------------------
\subsection{Connecting True Error to Training Error}
%-------------------------------------------------------------------

Using Stein's Lemma (with some technical steps):
\[
  \frac{1}{n}\sum_{i=1}^{n}\varepsilon_i(\hat{f}(x_i) - f(x_i))
  = \frac{\sigma^2}{n}\sum_{i=1}^{n}\frac{\partial\hat{f}(x_i)}{\partial y_i}
\]

\paragraph{What is $\frac{\partial\hat{f}(x_i)}{\partial y_i}$?} It measures how much the prediction $\hat{f}(x_i)$ changes when the observation $y_i$ changes by a small amount.

\begin{itemize}
  \item \textbf{Complex model}: A small change in $y_i$ causes a large change in $\hat{f}$ (it fits tightly to training data) $\Rightarrow$ large $\frac{\partial\hat{f}}{\partial y_i}$
  \item \textbf{Simple model}: Less sensitive to individual data points $\Rightarrow$ small $\frac{\partial\hat{f}}{\partial y_i}$
\end{itemize}

\paragraph{True error expression:}
\[
  \boxed{
    \text{True error}
    = \text{Empirical train error}
    + \underbrace{\frac{2\sigma^2}{n}\sum_{i=1}^{n}\frac{\partial\hat{f}(x_i)}{\partial y_i}}_{\Reg(\text{model complexity})}
  }
\]

\begin{keybox}
This is the \textbf{mathematical basis for all regularisation methods}. The term $\Reg(\theta) = \frac{\sigma^2}{n}\sum_{i}\frac{\partial\hat{f}(x_i)}{\partial y_i}$ measures model complexity. During training, instead of minimising $\Loss_{\text{train}}(\theta)$ alone, we should minimise:
\[
  \min_{\theta}\;\bigl[\Loss_{\text{train}}(\theta) + \Reg(\theta)\bigr]
\]
where $\Reg(\theta)$ is high for complex models and small for simple models.
\end{keybox}

%-------------------------------------------------------------------
\subsection{All Regularisation Methods as Instances}
%-------------------------------------------------------------------

All the following methods are instances of the above framework (they all approximate or bound $\Reg(\theta)$):
\begin{enumerate}
  \item L2 regularisation (weight decay)
  \item Dataset augmentation (empirical method)
  \item Parameter sharing and tying
  \item Adding noise to inputs
  \item Adding noise to outputs (label smoothing)
  \item Early stopping
  \item Ensemble methods
  \item Dropout
\end{enumerate}

%=====================================================================
\section{L2 Regularisation in Depth}
%=====================================================================

\subsection{Overview}

For L2 regularisation:
\[
  \tilde{\Loss}(w)
  = \Loss(w) + \frac{\alpha}{2}\norm{w}^2
  = \underbrace{\Loss(w)}_{\text{empirical train error}}
  + \underbrace{\frac{\alpha}{2}\norm{w}^2}_{\Reg(\theta)}
\]

\textbf{Why does it relate to model complexity?} Reducing weights $\Rightarrow$ reducing freedom in the model $\Rightarrow$ lower sensitivity $\Rightarrow$ lower $\frac{\partial\hat{f}}{\partial y_i}$.

\paragraph{SGD update with L2:}
\[
  w_{t+1} = w_t - \eta\grad\Loss(w_t) - \eta\alpha w_t
  = (1-\eta\alpha)w_t - \eta\grad\Loss(w_t)
\]

The L2 penalty ``shrinks'' weights toward zero at every step -- often called \textbf{weight decay}.

\begin{notebox}
Implementation is trivial: just add $\alpha w$ to every gradient $\frac{\partial\Loss}{\partial w}$. Works for any architecture.
\end{notebox}

%=====================================================================
\section{Dataset Augmentation}
%=====================================================================

\subsection{Core Idea}

Minimising the empirical training error allows the model to perfectly memorise training data. To avoid this, we \textbf{augment} the training data using domain knowledge.

\paragraph{Key principle:} Certain transformations to an input do not change its label. We exploit this to create additional training examples:
\begin{itemize}
  \item Rotation (e.g.\ by $5^\circ$)
  \item Horizontal/vertical shift
  \item Blurring
  \item Changing some pixels
  \item Horizontal flip
\end{itemize}

\paragraph{Why it helps:}
\begin{itemize}
  \item More data $=$ better learning (reduces overfitting)
  \item Works well for image classification/object recognition
  \item Also shown to work well for speech
\end{itemize}

\paragraph{Limitations:}
\begin{itemize}
  \item For some tasks it may not be clear how to generate augmented data (e.g.\ medical EHR data, detecting ``warm-hearted people'')
  \item All augmentation must use domain knowledge about what transformations preserve the label
\end{itemize}

\begin{notebox}
All augmentation is done with \textbf{domain knowledge} -- you must know which transformations are valid. This is a \textbf{supervised} form of regularisation.
\end{notebox}

%=====================================================================
\section{Parameter Sharing and Tying}
%=====================================================================

\begin{itemize}
  \item \textbf{Parameter Sharing} (used in CNNs):
    \begin{itemize}
      \item Same filter applied at different positions of the image
      \item Or: same weight matrix acts on different input neurons
      \item Uses the same parameter in multiple computations
      \item Drastically reduces number of parameters
    \end{itemize}

  \item \textbf{Parameter Tying} (used in AEs):
    \begin{itemize}
      \item Encoder and decoder weights tied: $W^* = W^\top$
      \item Imposes an equality constraint between parameters
      \item Reduces parameters from $2d\times n$ to $d\times n$
      \item Tying is a specific form of sharing that reduces free parameters
    \end{itemize}
\end{itemize}

\begin{notebox}
Both parameter sharing and tying relate to model complexity $\Reg(\theta)$ via Stein's Lemma: fewer effective parameters $\Rightarrow$ less sensitivity to data changes $\Rightarrow$ lower model complexity penalty.
\end{notebox}

%=====================================================================
\section{Adding Noise to Inputs}
%=====================================================================

\subsection{Equivalence to L2 Regularisation}

Adding Gaussian noise to the input is \textbf{equivalent to weight decay (L2 regularisation)} for simple input-output neural networks.

\paragraph{Formal derivation:} Let $\tilde{x}_i = x_i + \varepsilon_i$ where $\varepsilon_i \sim \mathcal{N}(0, \sigma^2)$.

For a linear model $\hat{y} = \sum_i w_i x_i = w^\top x$:
\[
  \tilde{y} = w^\top\tilde{x}_i = w^\top x_i + \sum_i w_i\varepsilon_i
  = \hat{y} + \sum_i w_i\varepsilon_i
\]

Now compute $\E[(\tilde{y}-y)^2]$:
\begin{align*}
  \E[(\tilde{y}-y)^2]
  &= \E\!\left[\left(\hat{y} + \sum_i w_i\varepsilon_i - y\right)^2\right]\\
  &= \E[(\hat{y}-y)^2]
    + 2\E\!\left[(\hat{y}-y)\sum_i w_i\varepsilon_i\right]
    + \E\!\left[\left(\sum_i w_i\varepsilon_i\right)^2\right]
\end{align*}

The cross term vanishes since $\varepsilon_i \perp (\hat{y}-y)$ (noise is independent of predictions):
\[
  \E\!\left[(\hat{y}-y)\sum_i w_i\varepsilon_i\right] = 0
\]

For the last term (since $\varepsilon_i$ are i.i.d.\ with $\E[\varepsilon_i^2]=\sigma^2$ and $\E[\varepsilon_i\varepsilon_j]=0$ for $i\neq j$):
\[
  \E\!\left[\left(\sum_i w_i\varepsilon_i\right)^2\right]
  = \sum_i w_i^2\E[\varepsilon_i^2] + \sum_{i\neq j}w_iw_j\underbrace{\E[\varepsilon_i\varepsilon_j]}_{=0}
  = \sigma^2\sum_i w_i^2
  = \sigma^2\norm{w}^2
\]

Therefore:
\[
  \boxed{
    \E[(\tilde{y}-y)^2]
    = \E[(\hat{y}-y)^2] + \sigma^2\norm{w}^2
  }
\]

This is exactly the \textbf{L2 regularisation} loss with $\lambda = \sigma^2$.

\begin{keybox}
Adding Gaussian noise to inputs = L2 regularisation on weights. The noise variance $\sigma^2$ plays the role of the regularisation parameter $\lambda$.
\end{keybox}

Adding noise to inputs can also be viewed as \textbf{data augmentation}.

%=====================================================================
\section{Adding Noise to Outputs (Label Smoothing)}
%=====================================================================

\subsection{Hard vs.\ Soft Targets}

For a classification problem with $K$ classes, the true distribution (hard targets / one-hot encoding):
\[
  p = [0, 0, \ldots, 1, \ldots, 0]
\]

The model estimates distribution $q$ from output layer.

\paragraph{Intuition:} Do not fully trust the labels -- they may be noisy. Instead use \textbf{soft targets}:
\[
  p = \left[\frac{\varepsilon}{K}, \ldots, 1-\varepsilon, \ldots, \frac{\varepsilon}{K}\right]
  \quad\text{(soft targets)}
\]
where $\varepsilon$ is a small positive constant representing label noise.

The cross-entropy loss with soft targets:
\[
  \min\;\sum_i p_i \log q_i
\]

\paragraph{Why does this work?} The model is not allowed to drive the training error to zero (because the soft targets have uncertainty). This prevents \textbf{overconfidence} and acts as regularisation.

\paragraph{The principle:} Any method of regularisation essentially says: ``\textbf{do not overfit under the training data}'' -- do not allow the true training error to go to zero.

%=====================================================================
\section{Early Stopping}
%=====================================================================

\subsection{Algorithm}

\begin{algorithm}
\caption{Early Stopping}
\begin{algorithmic}[1]
\State Track the validation error during training
\State Have a patience parameter $P$ (e.g.\ $P = 5$)
\State Store the model at every epoch
\If{at step $k$, no improvement in validation error for previous $P$ steps}
  \State \textbf{Stop training}
  \State \textbf{Return} model stored at step $k - P$
\EndIf
\end{algorithmic}
\end{algorithm}

\paragraph{Example (Patience = 5):}
\begin{center}
\begin{tabular}{cc}
\toprule
Epoch & Validation Error\\
\midrule
10 & 0.25\\
11 & 0.23 $\leftarrow$ best \\
12 & 0.24\\
13 & 0.25\\
14 & 0.25\\
15 & 0.24\\
16 & 0.26\\
\bottomrule
\end{tabular}
\end{center}

Stop at epoch 16 (no improvement for $P=5$ steps after epoch 11), return model from epoch 11.

%-------------------------------------------------------------------
\subsection{Early Stopping as Approximation to L2 Regularisation}
%-------------------------------------------------------------------

\begin{theorembox}[title={Early Stopping $\approx$ L2 Regularisation}]
Starting from $w_0 = 0$, with the SGD update rule and the quadratic approximation $\grad\Loss(w) \approx H(w-w^*)$:

After $T$ steps:
\[
  w_T = Q(I - (I-\eta\Lambda)^T)Q^\top w^*
\]

Compare with the L2 optimal solution:
\[
  \tilde{w} = Q(\Lambda+\alpha I)^{-1}\Lambda\, Q^\top w^*
\]

They match ($w_T = \tilde{w}$) when we choose $\eta, T, \alpha$ such that:
\[
  (I - \eta\Lambda)^T = (\Lambda + \alpha I)^{-1}\alpha
\]
\end{theorembox}

\paragraph{Interpretation:}
\begin{itemize}
  \item Early stopping only allows updates to parameters in directions that are \textbf{important for the loss}
  \item If a parameter direction has large eigenvalue $\lambda_i$: gradient is large, parameter grows $\Rightarrow$ direction is important, retained
  \item If a parameter direction has small eigenvalue $\lambda_i$: gradient is small, update is small, stays near zero $\Rightarrow$ direction is unimportant, effectively removed
  \item We do a limited number of $T$ steps
\end{itemize}

\begin{keybox}
Early stopping is the \textbf{most widely used} form of regularisation. It:
\begin{itemize}
  \item Is very effective
  \item Can be combined with other forms of regularisation (e.g.\ L2)
  \item Has a mathematical correspondence to L2 regularisation via spectral analysis
\end{itemize}
\end{keybox}

%=====================================================================
\section{Ensemble Methods and Bagging}
%=====================================================================

\subsection{Ensemble Methods}

\textbf{Combine the outputs of different models} to reduce generalisation error.

Models in an ensemble can be:
\begin{itemize}
  \item Different instances of the same classifier with different hyperparameters
  \item Models using different features
  \item Models trained on different samples of training data
\end{itemize}

%-------------------------------------------------------------------
\subsection{Bagging: Bootstrap Aggregating}
%-------------------------------------------------------------------

Form an ensemble of different instances of the \emph{same} classifier. From a given dataset, construct multiple training sets by \textbf{sampling with replacement}: $\mathcal{T}_1, \mathcal{T}_2, \ldots, \mathcal{T}_K$.

\textbf{When does bagging work?} Each classifier makes different errors at different positions. Combining reduces error. We want errors to be as \textbf{independent} as possible.

%-------------------------------------------------------------------
\subsection{Bagging MSE Analysis}
%-------------------------------------------------------------------

Consider $K$ models, each making error $\varepsilon_i$ on a test example. Let $\varepsilon_i$ be drawn from a zero-mean multivariate normal with:
\begin{itemize}
  \item Variance: $\E[\varepsilon_i^2] = V$
  \item Covariance: $\E[\varepsilon_i\varepsilon_j] = C$ for $i \neq j$
\end{itemize}

The expected squared error of the ensemble average $\frac{1}{K}\sum_i\varepsilon_i$:

\begin{align*}
  \text{MSE}
  &= \E\!\left[\left(\frac{1}{K}\sum_i\varepsilon_i\right)^2\right]\\
  &= \frac{1}{K^2}\E\!\left[\sum_i\varepsilon_i^2 + \sum_{i\neq j}\varepsilon_i\varepsilon_j\right]\\
  &= \frac{1}{K^2}\left[K\cdot\E[\varepsilon_i^2] + K(K-1)\cdot\E[\varepsilon_i\varepsilon_j]\right]\\
  &= \frac{1}{K^2}\left[KV + K(K-1)C\right]\\
  &= \frac{V}{K} + \frac{K-1}{K}C
\end{align*}

\begin{theorembox}[title={Bagging MSE}]
\[
  \text{MSE}_{\text{ensemble}} = \frac{V}{K} + \frac{K-1}{K}C
\]
\begin{itemize}
  \item If errors are \textbf{independent}: $C \approx 0$ $\Rightarrow$ $\text{MSE} \approx V/K$ (reduces by factor $K$!)
  \item If errors are \textbf{perfectly correlated}: $C = V$ $\Rightarrow$ $\text{MSE} = V$ (no improvement from ensembling)
\end{itemize}
\end{theorembox}

%=====================================================================
\section{Dropout}
%=====================================================================

\subsection{Motivation}

Training multiple large neural networks for an ensemble is \textbf{prohibitively expensive}. At test time, combining several models in real-time applications is infeasible.

\textbf{Dropout} is a technique that addresses both issues:
\begin{itemize}
  \item Effectively trains several neural networks \textbf{without significant computational overhead}
  \item Gives an efficient approximate way of combining \textbf{exponentially many} ($2^n$) different neural networks
\end{itemize}

%-------------------------------------------------------------------
\subsection{The Dropout Mechanism}
%-------------------------------------------------------------------

\textbf{Dropout} refers to \textbf{dropping out units} (neurons):
\begin{itemize}
  \item Temporarily remove a node and its incoming/outgoing connections
  \item Results in a \textbf{thinned network}
  \item Each node is retained with a fixed probability (typically $p = 0.5$ for hidden nodes, $p = 0.8$ for visible/input nodes)
\end{itemize}

Formally, for each neuron $j$:
\[
  \tilde{h}_j = r_j \cdot h_j,\quad r_j \sim \text{Bernoulli}(p)
\]

%-------------------------------------------------------------------
\subsection{Training Procedure}
%-------------------------------------------------------------------

\begin{enumerate}
  \item Initialise all parameters (weights) of the network
  \item For the \textbf{first training instance/mini-batch}: apply dropout $\Rightarrow$ thinned network $T_1$
  \item Compute loss on $T_1$ and backpropagate to \textbf{active weights only}
  \item For the \textbf{second training instance/mini-batch}: apply dropout again $\Rightarrow$ \textbf{different} thinned network $T_2$
  \item Compute loss on $T_2$ and backpropagate to active weights
  \item Repeat for all training instances
\end{enumerate}

\paragraph{Key observation:} With $n$ nodes, there are $2^n$ possible thinned networks. This is prohibitively large, but two tricks make it tractable:
\begin{enumerate}
  \item \textbf{Share weights} across all $2^n$ networks
  \item \textbf{Sample a different network} for each training instance
\end{enumerate}

%-------------------------------------------------------------------
\subsection{Why Dropout Works: Preventing Co-Adaptation}
%-------------------------------------------------------------------

\begin{keybox}
\textbf{Preventing Co-Adaptation:} Dropout prevents hidden units from co-adapting (relying too much on each other). A hidden unit cannot rely on other units, as they may be dropped at any time. Each unit must learn to be \textbf{more robust} to random dropouts.
\end{keybox}

\paragraph{Concrete example -- face detection:}
\begin{itemize}
  \item Suppose neuron $h_i$ learns to detect a face by firing when it detects a nose (most distinctive feature)
  \item All other neurons $h_j$ become lazy and rely on $h_i$
  \item If $h_i$ is dropped, other neurons must take responsibility -- they learn other features (eyes, ears, etc.) for redundancy and robustness
  \item This forces the model to learn truly general features
\end{itemize}

%-------------------------------------------------------------------
\subsection{Test Time: Weight Scaling}
%-------------------------------------------------------------------

At test time, aggregating outputs of $2^n$ thinned networks is impossible. Instead:
\begin{enumerate}
  \item Use the \textbf{full network} (all neurons active)
  \item Scale the output of each node by $p$ (the fraction of time it was active during training)
\end{enumerate}

\[
  w_{\text{test}} = p \cdot w_{\text{train}}
\]

This approximates averaging over all $2^n$ thinned networks.

%-------------------------------------------------------------------
\subsection{Parameter Sharing Interpretation}
%-------------------------------------------------------------------

All $2^n$ thinned networks \textbf{share parameters}. Over time:
\begin{itemize}
  \item If a weight was active in both training instances 1 and 2: it received two updates
  \item If a weight was active in only one: it received one update
  \item Each thinned network is trained rarely (or even never), but \textbf{parameter sharing ensures no model has untrained parameters}
\end{itemize}

\begin{notebox}
Dropout essentially applies a \textbf{masking noise} to the hidden units (analogous to denoising autoencoders, which apply masking noise to inputs).
\end{notebox}

%=====================================================================
\section{Discrete Cosine Transform (DCT) and JPEG Compression}
%=====================================================================

\subsection{Introduction and Objectives}

\begin{itemize}
  \item Represent an image as a \textbf{weighted sum of basis images} (cosine components)
  \item \textbf{Compression}: Set some DCT coefficients to zero (those the eye cannot perceive)
  \item This is the basis of \textbf{JPEG} compression
\end{itemize}

\begin{notebox}
Both DCT/JPEG and autoencoders aim for compact representations:
\begin{itemize}
  \item \textbf{DCT}: Fixed cosine basis (no learning, mathematically defined)
  \item \textbf{AE}: Learned basis (data-adaptive, can capture domain-specific structure)
\end{itemize}
\end{notebox}

%-------------------------------------------------------------------
\subsection{2D DCT Formula}
%-------------------------------------------------------------------

The 2D discrete cosine transform for an $N\times N$ image $f(x,y)$:
\[
  F(u,v) = \frac{2}{N}C(u)C(v)
    \sum_{x=0}^{N-1}\sum_{y=0}^{N-1}
      f(x,y)
      \cos\!\left(\frac{(2x+1)u\pi}{2N}\right)
      \cos\!\left(\frac{(2y+1)v\pi}{2N}\right)
\]

where:
\begin{itemize}
  \item $f(x,y)$: pixel value at spatial position $(x,y)$
  \item $F(u,v)$: DCT coefficient at frequency $(u,v)$
  \item $N$: size of the image block
  \item $x, y$: spatial coordinates
  \item $u, v$: frequency coordinates
  \item $C(k) = \begin{cases} 1/\sqrt{2} & k = 0\\ 1 & k > 0 \end{cases}$
\end{itemize}

%-------------------------------------------------------------------
\subsection{Basis Image Interpretation}
%-------------------------------------------------------------------

Each 2D DCT basis image $C_{k,\ell}(m,n)$ is a product of:
\begin{itemize}
  \item A vertical frequency $\frac{k}{2N}$ cosine
  \item A horizontal frequency $\frac{\ell}{2N}$ cosine
\end{itemize}

The image is a weighted sum of basis images:
\[
  f(x,y) = \sum_{k=0}^{N-1}\sum_{\ell=0}^{N-1}b_{k\ell}\,C_{k,\ell}(x,y)
\]

where $b_{k\ell}$ are the \textbf{DCT coefficients} (weights).

\paragraph{Key example -- Horizontal Edge:}
An image with a horizontal edge has energy concentrated in the \emph{first column} ($\ell=0$) of the DCT matrix, because the image only varies in the vertical direction and has no horizontal frequency content.

%-------------------------------------------------------------------
\subsection{Numerical Example: $2\times 2$ DCT}
%-------------------------------------------------------------------

\begin{examplebox}
Compute the 2D DCT for:
\[
  f(x,y) = \begin{bmatrix} 1 & 2\\ 3 & 4 \end{bmatrix},
  \quad N = 2
\]

\textbf{Computing $F(0,0)$} ($u=0, v=0$):
\begin{align*}
  F(0,0)
  &= C(0)C(0)\sum_{x=0}^{1}\sum_{y=0}^{1}f(x,y)
    \cos\!\left(\frac{(2x+1)\cdot0\cdot\pi}{4}\right)
    \cos\!\left(\frac{(2y+1)\cdot0\cdot\pi}{4}\right)\\
  &= \frac{1}{\sqrt{2}}\cdot\frac{1}{\sqrt{2}}(1+2+3+4) = \frac{1}{2}\cdot10 = 5
\end{align*}

\textbf{Computing $F(0,1)$} ($u=0, v=1$):
Using $\cos(0)=1$, $\cos(\pi/4) \approx 0.707$, $\cos(3\pi/4) \approx -0.707$:
\[
  F(0,1) = \frac{1}{\sqrt{2}}\cdot 1\cdot
  \left[1\cdot1\cdot0.707 + 2\cdot1\cdot(-0.707) + 3\cdot1\cdot0.707 + 4\cdot1\cdot(-0.707)\right]
  \approx -1
\]

\textbf{Full DCT result:}
\[
  F = \begin{bmatrix} 5 & -1\\ -2 & 0 \end{bmatrix}
\]
\end{examplebox}

%-------------------------------------------------------------------
\subsection{JPEG Compression: Step-by-Step Pipeline}
%-------------------------------------------------------------------

\subsubsection{Step 1: Colour Space Conversion (RGB $\to$ YCbCr)}

Convert from RGB $\in [0,255]$ to YCbCr:
\begin{itemize}
  \item $Y$: luminance (brightness) -- humans are more receptive to this
  \item $C_b$: blue chrominance
  \item $C_r$: red chrominance -- humans perceive colour changes less sensitively
\end{itemize}
This conversion is \textbf{reversible} (no data loss yet).

\subsubsection{Step 2: Chrominance Downsampling}

Remove considerable colour data since humans perceive colour less finely:
\[
  C_b, C_r \to (2\times2)\text{ average}
\]
Each $2\times2$ block of $C_b$ and $C_r$ is replaced by a single average value.

\paragraph{Storage comparison:}
\begin{itemize}
  \item Before: $1 + 1 + 1 = 3$ values per pixel
  \item After: $1 + \frac{1}{4} + \frac{1}{4} = 1.5$ values per pixel
\end{itemize}
This gives a \textbf{2$\times$ compression} before any DCT!

\subsubsection{Step 3: Break up into $8\times8$ Blocks}

Break the $M\times N$ image into non-overlapping $8\times8$ blocks.

\subsubsection{Step 4: Take 2D DCT of Each Block}

Compute the DCT of each $8\times8$ block to obtain 64 DCT coefficients. The DCT \textbf{concentrates energy} in the low-frequency coefficients (top-left of the $8\times8$ matrix).

\subsubsection{Step 5: Quantisation}

The \textbf{quantisation} step is where compression actually happens:
\[
  \text{Quantised coefficient} = \mathrm{round}\!\left(\frac{\text{DCT coefficient}}{\text{Quantisation table entry}}\right)
\]

The \textbf{quantisation table}:
\begin{itemize}
  \item Lower values in top-left: more distinct patterns kept (low frequency -- important)
  \item Higher values in bottom-right: high-frequency detail rounded to 0 (eyes can't perceive it anyway)
\end{itemize}

After quantisation, many high-frequency coefficients become zero.

\subsubsection{Step 6: Run-Length and Huffman Encoding}

Quantised values are read in a zigzag order (concentrates zeros together):
\[
  140, -14, -60, 14, -7, \ldots, -2, 4(\times2), -1, 0, 2, -1, -3, \ldots
\]

\textbf{Run-Length Encoding (RLE)}: Encode as value, then number of zeros: e.g.\ $-1, 0[\times53]$.

\textbf{Huffman Encoding}: Variable-length codes for common values (shorter codes for more frequent values).

\subsubsection{JPEG Decoding}
\begin{enumerate}
  \item Read bits for each $8\times8$ block
  \item Convert bits to DCT coefficients
  \item Take inverse DCT to find the $8\times8$ image block
  \item Assemble blocks together
\end{enumerate}

%-------------------------------------------------------------------
\subsection{Why JPEG Works: Human Vision Properties}
%-------------------------------------------------------------------

\paragraph{Light receptor cells:}
\begin{itemize}
  \item \textbf{Rods} ($\sim$100 million): brightness sensitive
  \item \textbf{Cones} ($\sim$6 million): colour sensitive
\end{itemize}

Human eyes are \textbf{good at seeing} edges (rock edges, image boundaries) but \textbf{not at detecting} high-frequency colour differences such as:
\begin{itemize}
  \item Individual blades of grass
  \item Individual leaves in a cluster
  \item Variations in shadows from tree leaves
\end{itemize}

JPEG exploits this: remove information that our eyes cannot perceive.

\paragraph{Impact of JPEG:}
\begin{itemize}
  \item 86\% of all images on the internet are JPEGs
  \item Video compression algorithms such as H.264 (AVC = Advanced Video Coding) use similar ideas
  \item Saves servers \textbf{petabytes of storage} and \textbf{billions of dollars} in cost reduction
\end{itemize}

%=====================================================================
\section{Recap: Training Deep Neural Networks}
%=====================================================================

\subsection{Single Neuron (Sigmoid)}

For a single neuron with sigmoid activation $f(x) = \sigma(w^\top x)$:
\[
  w \leftarrow w - \eta\nabla w
  \quad\text{where}\quad
  \nabla w = (\hat{f}(x) - y)\cdot f(x)\cdot(1-f(x))\cdot x
\]

\subsection{Wider Network (Multiple Inputs)}

For a wider network with inputs $x_1, x_2, x_3$:
\[
  w_j \leftarrow w_j - \eta\frac{\partial\Loss}{\partial w_j},
  \quad j = 1, 2, 3
\]

\subsection{Deeper Networks: Chain Rule}

For a network with layer $l$, activations $h_l = \sigma(a_l)$, $a_l = w_l h_{l-1}$:
\[
  \frac{\partial\Loss}{\partial w_l}
  = \frac{\partial\Loss}{\partial h_l}
    \cdot \frac{\partial h_l}{\partial a_l}
    \cdot \frac{\partial a_l}{\partial w_l}
  = \delta_l \cdot h_{l-1}
\]

\begin{keybox}
\textbf{Key observation:} $\nabla w_l$ is \textbf{proportional to the corresponding input} $h_{l-1}$.

The gradient $\nabla w_i \propto h_{i-1}$ (the input to that layer). This is why backpropagation is elegant -- gradients flow backward while activations flow forward.
\end{keybox}

\subsection{Wide and Deep Networks}

For a network with $L$ layers and multiple nodes per layer:
\[
  \frac{\partial\Loss}{\partial w_l}
  = \text{sum of gradients over all paths from loss to } w_l \text{ (chain rule across multiple paths)}
\]

This is the essence of backpropagation -- it efficiently computes all these path sums.

\paragraph{Historical note:} Backpropagation was popularised by Rumelhart et al.\ in 1986. It did not catch on at that time due to insufficient computational power and memory. Until 2006, training very deep networks was extremely difficult and convergence was not guaranteed.

%=====================================================================
\section{Complete Summary}
%=====================================================================

\subsection{Autoencoders}

\begin{itemize}
  \item \textbf{Under-complete}: $\dim(\vh) < \dim(\vx)$ -- forces compression
  \item \textbf{Over-complete}: $\dim(\vh) \geq \dim(\vx)$ -- needs regularisation
  \item \textbf{Loss}: MSE (real inputs) or cross-entropy (binary inputs)
  \item \textbf{PCA equivalence}: linear encoder + linear decoder + MSE + normalised inputs
    $\Rightarrow H = XV_{\cdot,\leq k}$, $W = V_{\cdot,\leq k}^\top$
\end{itemize}

\subsection{Regularisation Methods for Autoencoders}

\begin{center}
\begin{tabular}{lp{6cm}}
\toprule
\textbf{Method} & \textbf{Mechanism} \\
\midrule
L2 / Weight Decay & $\Reg(\theta) = \lambda\norm{\theta}^2$; shrink weights \\
Weight Tying & $W^* = W^\top$; halves parameters \\
Denoising & Corrupt $\tilde{x}\sim P(\tilde{x}|x)$; must reconstruct original \\
Sparse (KL) & $\hat{\rho}_\ell = \rho$; neurons inactive most of time \\
Contractive & $\Reg = \norm{J_f(\vh)}_F^2$; penalise sensitivity \\
\bottomrule
\end{tabular}
\end{center}

\subsection{Bias--Variance}

\[
  \E[(y-\hat{f})^2] = \underbrace{(\Bias)^2}_{\text{underfitting}} + \underbrace{\Var}_{\text{overfitting}} + \underbrace{\sigma^2}_{\text{irreducible}}
\]

\[
  \text{True error} = \text{Train error} + \underbrace{\frac{2\sigma^2}{n}\sum_i\frac{\partial\hat{f}(x_i)}{\partial y_i}}_{\Reg(\text{model complexity})}
\]

Always use a \textbf{validation set} (independent of training set) to estimate error.

\subsection{General Regularisation Methods}

All the following methods reduce $\Reg(\text{model complexity})$:

\begin{center}
\begin{tabular}{lp{7cm}}
\toprule
\textbf{Method} & \textbf{Mechanism} \\
\midrule
L2 regularisation & Penalise large weights \\
Dataset augmentation & More (transformed) training data \\
Parameter sharing/tying & Fewer free parameters \\
Noise to inputs & Equivalent to L2 regularisation \\
Noise to outputs (label smoothing) & Don't trust labels fully \\
Early stopping & Stop before overfitting \\
Ensemble/Bagging & Average independent models \\
Dropout & Random thinned networks \\
\bottomrule
\end{tabular}
\end{center}

\subsection{Key Equations at a Glance}

\paragraph{Autoencoder training:}
\[
  \Loss = \frac{1}{m}\norm{X - \hat{X}}_F^2,\quad
  \hat{X} = f(g(X)),\quad
  \vh = g(W\vx + \vb),\quad
  \hatx = f(W^*\vh + \vc)
\]

\paragraph{PCA--AE Equivalence:}
\[
  W = V_{\cdot,\leq k}^\top,\quad
  H = XV_{\cdot,\leq k},\quad
  X = U\Sigma V^\top
\]

\paragraph{Bias--Variance:}
\[
  \E\!\left[(y-\hat{f}(x))^2\right]
  = (f(x) - \E[\hat{f}(x)])^2
  + \E[(\hat{f}(x) - \E[\hat{f}(x)])^2]
  + \sigma^2
\]

\paragraph{L2 regularised solution:}
\[
  \tilde{w} = Q(\Lambda + \alpha I)^{-1}\Lambda Q^\top w^*
  = Q\,D\,Q^\top w^*
\]
where $D_{ii} = \frac{\lambda_i}{\lambda_i + \alpha}$ (shrinks directions with small eigenvalues).

\paragraph{Dropout at test time:}
\[
  w_{\text{test}} = p \cdot w_{\text{train}}
\]

\paragraph{Bagging MSE:}
\[
  \text{MSE}_{\text{ensemble}}
  = \frac{V}{K} + \frac{K-1}{K}C
  \;\xrightarrow{C=0}\;
  \frac{V}{K}
  \;\xrightarrow{C=V}\;
  V
\]

\begin{center}
\Large\textit{``Regularisation: any method that reduces test error at the expense of training error.''}
\end{center}

%=====================================================================
\appendix
%=====================================================================

\section{Appendix A: Denoising AE -- Gaussian Noise Equivalence to L2}

\begin{theorembox}[title={Full Derivation}]
Let $\tilde{x}_i = x_i + \varepsilon_i$, $\varepsilon_i \sim \mathcal{N}(0, \sigma^2)$.

Original model: $\hat{y} = \sum_i w_i x_i$.
Noisy model: $\tilde{y} = \sum_i w_i\tilde{x}_i = \hat{y} + \sum_i w_i\varepsilon_i$.

\begin{align*}
  \E[(\tilde{y}-y)^2]
  &= \E\!\left[(\hat{y} + \sum_i w_i\varepsilon_i - y)^2\right]\\
  &= \E[(\hat{y}-y)^2]
    + 2\E\!\left[(\hat{y}-y)\sum_i w_i\varepsilon_i\right]
    + \E\!\left[\left(\sum_i w_i\varepsilon_i\right)^2\right]
\end{align*}

Cross term vanishes: $\varepsilon_i \perp (\hat{y}-y)$ since noise is independent of predictions:
\[
  2\E\!\left[(\hat{y}-y)\sum_i w_i\varepsilon_i\right] = 0
\]

Last term (i.i.d.\ $\varepsilon_i$, $\E[\varepsilon_i\varepsilon_j] = 0$ for $i\neq j$, $\E[\varepsilon_i^2]=\sigma^2$):
\[
  \E\!\left[\left(\sum_i w_i\varepsilon_i\right)^2\right]
  = \sum_i w_i^2\E[\varepsilon_i^2] = \sigma^2\norm{w}^2
\]

Therefore:
\[
  \E[(\tilde{y}-y)^2] = \E[(\hat{y}-y)^2] + \sigma^2\norm{w}^2
\]

This is exactly L2 regularisation with $\lambda = \sigma^2$. \hfill$\square$
\end{theorembox}

%-------------------------------------------------------------------
\section{Appendix B: Sparse AE -- KL Divergence Details}
%-------------------------------------------------------------------

\begin{theorembox}[title={KL Divergence Between Bernoullis}]
For $\mathrm{Bern}(\rho)$ vs $\mathrm{Bern}(\hat{\rho}_\ell)$:
\[
  \mathrm{KL}(\rho\|\hat{\rho}_\ell)
  = \rho\log\frac{\rho}{\hat{\rho}_\ell} + (1-\rho)\log\frac{1-\rho}{1-\hat{\rho}_\ell}
  \geq 0
\]
with equality iff $\hat{\rho}_\ell = \rho$.

Gradient:
\[
  \frac{\partial\Reg}{\partial\hat{\rho}_\ell}
  = -\frac{\rho}{\hat{\rho}_\ell} + \frac{1-\rho}{1-\hat{\rho}_\ell}
\]

Combined weight update:
\[
  \frac{\partial\tilde{\Loss}}{\partial w}
  = \frac{\partial\Loss}{\partial w}
  + \left(-\frac{\rho}{\hat{\rho}_\ell} + \frac{1-\rho}{1-\hat{\rho}_\ell}\right)
    \cdot \vx_i\bigl(g'(w^\top\vx_i+\vb)\bigr)^\top
\]
\end{theorembox}

%-------------------------------------------------------------------
\section{Appendix C: Contractive AE -- Jacobian Norm for Sigmoid Units}
%-------------------------------------------------------------------

For sigmoid hidden units $h_\ell = \sigma(w_\ell^\top\vx + b_\ell)$:
\[
  \frac{\partial h_\ell}{\partial x_j}
  = \underbrace{h_\ell(1-h_\ell)}_{\sigma'\text{-derivative}}
    \cdot\, w_{\ell j}
\]

Therefore:
\[
  \norm{J_f(\vh)}_F^2
  = \sum_{\ell=1}^{k}\sum_{j=1}^{n}
    \left(\frac{\partial h_\ell}{\partial x_j}\right)^2
  = \sum_{\ell=1}^{k}(h_\ell(1-h_\ell))^2\sum_{j=1}^{n}w_{\ell j}^2
  = \sum_{\ell=1}^{k}(h_\ell(1-h_\ell))^2\norm{w_\ell}^2
\]

Combined objective:
\[
  \min_\theta\;\Loss(\theta) + \lambda\sum_\ell(h_\ell(1-h_\ell))^2\norm{w_\ell}^2
\]

Trade-off:
\begin{itemize}
  \item $\Loss(\theta)$: want $\vh$ to capture all variations in data
  \item $\norm{J_f(\vh)}_F^2$: want $\vh$ to be insensitive to variations in $\vx$
  \item Together: $\vh$ only responds to the most important (high-variance) directions
\end{itemize}

%-------------------------------------------------------------------
\section{Appendix D: Stein's Lemma -- Formal Statement and Application}
%-------------------------------------------------------------------

\begin{theorembox}[title={Stein's Lemma (Formal)}]
If $X \sim \mathcal{N}(\mu, \sigma^2)$ and $g$ is a differentiable function with $\E[|g'(X)|] < \infty$, then:
\[
  \E[(X-\mu)g(X)] = \sigma^2\,\E[g'(X)]
\]
\end{theorembox}

\paragraph{Application to true error estimation:} Using Stein's Lemma:
\[
  \frac{1}{n}\sum_{i=1}^{n}\varepsilon_i(\hat{f}(x_i) - f(x_i))
  = \frac{\sigma^2}{n}\sum_{i=1}^{n}\frac{\partial\hat{f}(x_i)}{\partial y_i}
\]

This shows:
\[
  \text{True error}
  = \underbrace{\text{Emp.\ train error}}_{\frac{1}{n}\sum(\hat{y}_i-y_i)^2}
  + \underbrace{\frac{2\sigma^2}{n}\sum_{i=1}^{n}\frac{\partial\hat{f}(x_i)}{\partial y_i}}_{\Reg(\theta)}
\]

Complex models have large $\frac{\partial\hat{f}(x_i)}{\partial y_i}$ (sensitive to data changes), while simple models have small values. Regularisation controls $\Reg(\theta)$.

%-------------------------------------------------------------------
\section{Appendix E: Bagging MSE -- Full Derivation}
%-------------------------------------------------------------------

Setup: $K$ models, each making error $\varepsilon_i$ on a test example, drawn from $\mathcal{N}_K(0, \Sigma)$ with $\Var(\varepsilon_i) = V$ and $\Cov(\varepsilon_i, \varepsilon_j) = C$ for $i \neq j$.

\begin{align*}
  \text{MSE}
  &= \E\!\left[\left(\frac{1}{K}\sum_i\varepsilon_i\right)^2\right]\\
  &= \frac{1}{K^2}\E\!\left[\sum_{i=j}\varepsilon_i^2 + \sum_{i\neq j}\varepsilon_i\varepsilon_j\right]\\
  &= \frac{1}{K^2}\left[K\cdot\E[\varepsilon_i^2] + K(K-1)\cdot\E[\varepsilon_i\varepsilon_j]\right]\\
  &= \frac{1}{K^2}[KV + K(K-1)C]\\
  &= \frac{V}{K} + \frac{K-1}{K}C
\end{align*}

When $C = 0$ (independent): MSE $= V/K$ (reduces by factor $K$).\\
When $C = V$ (perfectly correlated): MSE $= V$ (no improvement from ensembling).

%-------------------------------------------------------------------
\section{Appendix F: Early Stopping -- L2 Correspondence (Full Derivation)}
%-------------------------------------------------------------------

\paragraph{Setup:} Starting from $w_0 = 0$, with update rule:
\[
  w_{t+1} = w_t - \eta\grad\Loss(w_t)
\]

Using the quadratic approximation $\grad\Loss(w) \approx H(w-w^*)$:
\[
  w_{t+1} = w_t - \eta H(w_t - w^*) = (I-\eta H)w_t + \eta Hw^*
\]

Using SVD of $H = Q\Lambda Q^\top$:
\[
  w_t = Q(I-(I-\eta\Lambda)^t)Q^\top w^*
\]

The scaling factor for direction $i$ at step $T$: $1 - (1-\eta\lambda_i)^T$.

\paragraph{Comparison with L2 optimal:}
L2 regularised optimal: $\tilde{w} = Q(\Lambda+\alpha I)^{-1}\Lambda Q^\top w^*$, scaling factor $\frac{\lambda_i}{\lambda_i+\alpha}$.

They match when:
\[
  1 - (1-\eta\lambda_i)^T = \frac{\lambda_i}{\lambda_i+\alpha}
\]
i.e.\ for appropriate $\eta, T, \alpha$.

\paragraph{Intuitive interpretation:}
\begin{itemize}
  \item Large $\lambda_i$: $(1-\eta\lambda_i)^T \approx 0$ $\Rightarrow$ scaling factor $\approx 1$ (direction fully updated, retained)
  \item Small $\lambda_i$: $(1-\eta\lambda_i)^T \approx 1$ $\Rightarrow$ scaling factor $\approx 0$ (direction not updated, effectively zero)
\end{itemize}

This is exactly what L2 does: retain important directions, shrink unimportant ones.

%-------------------------------------------------------------------
\section{Appendix G: Backpropagation Through AE (Cross-Entropy + Sigmoid)}
%-------------------------------------------------------------------

For cross-entropy loss with sigmoid decoder $h_{2j} = \sigma(a_{2j})$:

\textbf{Gradient at output:}
\[
  \frac{\partial\Loss}{\partial a_{2j}}
  = \hat{x}_{ij} - x_{ij}
\]
(Clean formula: just the prediction error!)

\textbf{Gradient for decoder weights:}
\[
  \frac{\partial\Loss}{\partial W^*}
  = \frac{\partial\Loss}{\partial\vh_2}\cdot\frac{\partial\vh_2}{\partial a_2}\cdot\frac{\partial a_2}{\partial W^*}
  = (\hatx_i - \vx_i)\cdot\vh^\top
\]

\textbf{Gradient for encoder weights:}
\[
  \frac{\partial\Loss}{\partial W}
  = \delta_e \cdot \vx_i^\top
\]
where $\delta_e$ is computed by chaining backward through the decoder.

The autoencoder is just a special case of a feed-forward network with the same input and target, trained using standard backpropagation.

%-------------------------------------------------------------------
\section{Appendix H: PCA--AE Equivalence -- Step-by-Step with Dimensions}
%-------------------------------------------------------------------

\paragraph{Dimension tracking:}
\begin{itemize}
  \item $X \in \R^{m\times n}$ (data matrix: $m$ samples, $n$ features)
  \item $W \in \R^{k\times n}$ (encoder: maps $n$-dim to $k$-dim)
  \item $H = XW^\top \in \R^{m\times k}$ (hidden representations)
  \item $W^* \in \R^{k\times n}$ (decoder: maps $k$-dim back to $n$-dim)
  \item $\hat{X} = HW^* \in \R^{m\times n}$ (reconstructions)
\end{itemize}

\paragraph{SVD of $X$:}
$X = U\Sigma V^\top$ where:
\begin{itemize}
  \item $U \in \R^{m\times m}$: left singular vectors, $UU^\top = U^\top U = I$
  \item $\Sigma \in \R^{m\times n}$: singular values (diagonal)
  \item $V \in \R^{n\times n}$: right singular vectors, $VV^\top = V^\top V = I$
\end{itemize}

\paragraph{Optimal $\hat{X}$ (Eckart-Young):}
$\hat{X} = U_{\cdot,\leq k}\Sigma_{k\times k}V_{\cdot,\leq k}^\top$

\paragraph{Solution: $W^* = V_{\cdot,\leq k}^\top \in \R^{k\times n}$}

\paragraph{Deriving encoder $W$:} Show $H = XV_{\cdot,\leq k}$ implies $W = V_{\cdot,\leq k}^\top \in \R^{k\times n}$ (since $H = WX$ requires $W = H X^+ = V_{\cdot,\leq k}^\top$).

\paragraph{PCA connection:} $V = $ right singular vectors of $X$ $=$ eigenvectors of $X^\top X$ (covariance matrix after normalisation). PCA projects onto eigenvectors of covariance. Hence $W = V_{\cdot,\leq k}^\top = P$ (PCA projection matrix). \hfill\textbf{QED}

\end{document}